% MIRROR COPY. The canonical source of this paper lives in
% ~/GitHub/fiber-collapsed-mcmc/paper/balanced_fiber_gibbs.tex
% (git-tracked, with the verification suite). Edit there and re-copy.
\documentclass[11pt]{article}

\usepackage{amssymb,amsmath,amsfonts,amsthm}
\usepackage{lmodern}
\usepackage{xcolor}
\usepackage[colorlinks=true]{hyperref}
\definecolor{cieblue}{RGB}{31, 119, 180}
\hypersetup{citecolor=cieblue, allcolors=cieblue}
\usepackage{enumerate}

\setlength{\textwidth}{15.5cm}
\setlength{\textheight}{22cm}
\setlength{\oddsidemargin}{0.25cm}
\setlength{\evensidemargin}{0.25cm}
\setlength{\topmargin}{0cm}

% ---------- status annotations (REMOVE BEFORE SUBMISSION) ----------
\definecolor{donecol}{RGB}{0,110,90}
\definecolor{misscol}{RGB}{170,20,20}
\definecolor{notecol}{RGB}{30,60,170}
\newcommand{\done}[1]{\medskip\noindent\textcolor{donecol}{\textbf{[DONE]} \emph{#1}}\medskip}
\newcommand{\missing}[1]{\medskip\noindent\textcolor{misscol}{\textbf{[MISSING]} \emph{#1}}\medskip}
\newcommand{\statusnote}[1]{\medskip\noindent\textcolor{notecol}{\textbf{[NOTE]} \emph{#1}}\medskip}
% -------------------------------------------------------------------

\newcommand{\Z}{\mathbb{Z}}
\newcommand{\R}{\mathbb{R}}
\newcommand{\E}{\mathbb{E}}
\newcommand{\PicO}[1]{\operatorname{Pic}^0(#1)}
\newcommand{\Pick}[1]{\operatorname{Pic}_k(#1)}
\newcommand{\im}{\operatorname{im}}
\newcommand{\tv}[1]{\left\lVert #1 \right\rVert_{\mathrm{TV}}}
\newcommand{\Tcal}{\mathcal{T}}
\newcommand{\Fcal}{\mathcal{F}}
\newcommand{\Bcal}{\mathcal{B}}

\newtheorem{theorem}{Theorem}[section]
\newtheorem{lemma}[theorem]{Lemma}
\newtheorem{proposition}[theorem]{Proposition}
\newtheorem{corollary}[theorem]{Corollary}
\newtheorem{conjecture}[theorem]{Conjecture}
\newtheorem{definition}[theorem]{Definition}
\newtheorem{remark}[theorem]{Remark}
\newtheorem{fact}[theorem]{Fact}

\title{A Multi-Sink Sandpile Torsor and Fiber-Collapsed Markov Chains\\
for Sampling Balanced Graph Partitions\\
{\large \textcolor{misscol}{WORKING DRAFT, July 2026 (see the Status appendix)}}}
\author{Alaittin K\i rt\i \c{s}o\u{g}lu\thanks{Department of Applied
Mathematics, Illinois Institute of Technology.}}
\date{July 2026}

\begin{document}
\maketitle

\begin{abstract}
Sampling connected, population-balanced $k$-partitions of a graph, the
mathematical core of political redistricting analysis, currently rests on
Markov chains with no known mixing-time bounds; rigorous guarantees exist only
for \emph{unconstrained} spanning-forest walks. We make three contributions
toward closing this gap. First, we develop the algebraic infrastructure of
balanced-forest chains. We prove that the lattice quotient
$\Pick{G} = \Z^{|V|-k}/\im \Delta^{(S)}$ attached to a root set $S$ of size $k$
is the sandpile group of the graph $G_S$ obtained by identifying $S$ to a
single sink, that spanning $k$-forests of $G$ rooted at $S$ are exactly the
spanning trees of $G_S$, and hence, by rotor-routing on $G_S$, that
$\Pick{G}$ acts freely and transitively on rooted $k$-forests of any graph.
This yields exact (Markov-chain-free) uniform sampling of rooted $k$-forests
and, via the quotient multigraph $Q_\xi$, an \emph{exact within-fiber
resampler} for the set of forests or trees compatible with a fixed partition
$\xi$. Second, we prove a \emph{fiber collapse lemma}. A chain that
alternates exact fiber resampling with an arbitrary reversible
partition-changing kernel mixes \emph{exactly} as fast as its projection
onto partition space; the within-fiber coordinate contributes nothing to
the mixing time. This eliminates, provably and by construction, one of the two
sources of slow mixing in all existing redistricting chains, and reduces the
open problem of balanced-chain mixing to a spectral gap estimate for a
projected chain on partitions alone. Third, on $2 \times n$ strips we
complete the program. We classify the balanced connected partitions of the strip
exactly for every $k$, by an interface coding with $\Theta_k(n^{k-1})$
states whose stationary law concentrates exponentially on narrow
interfaces. We prove that the projected window-resampling chain is
irreducible for every $k$, that every merged window splits with
probability at least $1/10$, and, by a two-phase coupling driven by the
concentration, that the full chain mixes. For every fixed $k$, the
fiber-Gibbs chain on balanced $k$-forests of the $2 \times n$ strip
mixes in $O(k \log k)$ steps, so an approximately uniform sample costs
near-linear time. To our knowledge this is the first polynomial
mixing-time bound for a Markov chain supported on balanced connected
partitions. Exact diagonalization shows that the projected spectral gap
is in fact $1/(k-1)$ asymptotically, independent of $n$. What remains
open on two-dimensional grids is stated as a precise conjecture.
\end{abstract}

\statusnote{This draft merges the ``Tier 2'' package (torsor infrastructure,
exact fiber resampling, irreducibility) with the ``Tier 1'' target (first
mixing bound for a balance-restricted chain). Everything labeled DONE has a
complete proof in this file or a precise citation. Everything labeled
MISSING is work to do; each MISSING box says exactly what is needed. The
ledger in Appendix~\ref{app:status} lists all of them in one table. Remove
all colored boxes, this note, and Appendix~\ref{app:status} before
submission.}

\tableofcontents

%═══════════════════════════════════════════════════════════════════════════════
\section{Introduction}
%═══════════════════════════════════════════════════════════════════════════════

Let $G = (V, E)$ be a connected graph with $n = |V|$ vertices and $m = |E|$
edges, and let $k \ge 2$. A \emph{connected $k$-partition} of $G$ is a
partition $\xi = \{\xi_1, \ldots, \xi_k\}$ of $V$ into $k$ nonempty classes
(\emph{districts}), each inducing a connected subgraph. A partition is
\emph{$\varepsilon$-balanced} if every district's size (more generally,
vertex-weighted population) lies within a factor $(1 \pm \varepsilon)$ of
$n/k$. The central algorithmic problem of quantitative redistricting analysis
is to sample from a prescribed distribution on balanced connected
$k$-partitions, most prominently the \emph{spanning tree distribution}
\[
  \pi_{\mathrm{sp}}(\xi) \;\propto\; \prod_{i=1}^{k} t(G[\xi_i]),
\]
where $t(H)$ denotes the number of spanning trees of $H$, which arises as the
projection of the uniform distribution on $k$-component spanning forests.
Chains in practical use (recombination~\cite{deford2021recom}, its
reversible forest variants~\cite{autry2020multiscale, charikar2023complexity},
the cycle walk~\cite{deford2025cycle}, the marked edge walk~\cite{mew}, and
the balanced up-down walk~\cite{akitaya2026bud}) are supported by
increasingly strong irreducibility theorems and extensive experiments, but
\emph{no polynomial mixing-time bound is known for any Markov chain whose
state space is the balanced set}, on any nontrivial graph family. The only
rigorous mixing results in the area concern \emph{unconstrained} forest
walks: the up-down walk on the rank-$(n-k)$ truncation of the graphic matroid
mixes in $O(m \log m)$ steps on every
graph~\cite{anari2021logconcave}, and Cannon, Pegden, and
Tucker-Foltz~\cite{cannon2024balanced} showed that on $m \times n$ grids the
balanced states occupy a $1/\mathrm{poly}$ fraction of the unconstrained
space, so that rejection from the unconstrained sampler runs in expected
polynomial time. Whether a chain \emph{living inside} the balanced set can be
proved to mix rapidly is open; it is stated as the outstanding problem
in~\cite{deford2025cycle, akitaya2026bud, cannon2024balanced}. This paper
resolves it for $2 \times n$ strips, for every fixed number of
districts.

Every chain on balanced forests suffers from two distinct sources of slow
mixing, and existing analyses cannot separate them:
\begin{enumerate}[(S1)]
  \item \textbf{within-fiber shuffling:} states carrying the same partition
    $\xi$ (differing in within-district trees, or in which boundary edges the
    spanning structure uses) must be mixed among themselves; empirically this
    consumes the majority of the steps of tree-based
    chains~\cite{mew, deford2025cycle};
  \item \textbf{partition-space geometry:} the projected walk on balanced
    partitions themselves must overcome the combinatorial geometry of the
    balanced set, the genuinely hard part, where torpidity examples
    exist~\cite{charikar2023complexity} and irreducibility can even
    fail~\cite{akitaya2026bud}.
\end{enumerate}

\paragraph{Contributions.} This paper removes (S1) \emph{exactly}, by
algebraic means, and isolates (S2) as a clean projected-chain problem.

\begin{enumerate}[(C1)]
  \item \textbf{A multi-sink sandpile torsor
    (Theorem~\ref{thm:torsor}).} For any root set
    $S = \{q_1, \ldots, q_k\} \subseteq V$, the lattice quotient
    $\Pick{G} := \Z^{n-k} / \im \Delta^{(S)}$ is canonically the sandpile
    (Picard) group of the multigraph $G_S$ obtained by identifying $S$ to a
    single sink; spanning $k$-forests of $G$ rooted at $S$ are exactly the
    spanning trees of $G_S$; and rotor-routing on $G_S$ makes the set of
    rooted $k$-forests a \emph{torsor} (free transitive action) for
    $\Pick{G}$, on every connected graph, with no planarity hypothesis.
    Consequences: exact uniform sampling of rooted $k$-forests by Wilson's
    algorithm on $G_S$ (no Markov chain, no mixing bound), a multi-sink
    Dhar burning algorithm, and polynomial-time navigation of forest space.
  \item \textbf{Exact fiber resampling
    (Theorem~\ref{thm:fiber-decomp}, Proposition~\ref{prop:resampler}).}
    The set of spanning trees compatible with a fixed partition $\xi$
    factorizes as
    $\prod_i \Tcal(G[\xi_i]) \times \Tcal(Q_\xi)$, where $Q_\xi$ is the
    quotient multigraph; both factors admit exact uniform samplers (Wilson
    within districts; Wilson or the $\PicO{Q_\xi}$-torsor on the quotient).
    This gives a polynomial-time \emph{exact} sampler of the conditional
    distribution given $\xi$, for both forest-valued and tree-valued state
    spaces.
  \item \textbf{The fiber collapse lemma (Lemma~\ref{lem:collapse},
    Theorem~\ref{thm:conditional}).} If a chain alternates the exact fiber
    resampler with an arbitrary $\pi$-reversible partition-changing kernel
    $P_M$, then the induced partition process is itself Markov, the joint
    chain's distance to stationarity \emph{equals} that of the projected
    chain, and the spectrum of the joint kernel is the spectrum of the
    projected kernel together with $0$. With exact fiber resampling,
    the mixing problem for balanced forests \emph{is} the mixing problem
    for balanced partitions, and (S1) is gone.
  \item \textbf{The first mixing bound for a balanced-partition chain
    (Section~\ref{sec:targets}).} On $2 \times n$ strips we classify the
    balanced partitions exactly for every $k$
    (Proposition~\ref{prop:strip-classify-k}), prove irreducibility of
    the projected window-resampling chain for every $k$
    (Corollary~\ref{cor:strip-irred}), and prove that the fiber-Gibbs
    chain mixes in $O(k \log k)$ steps
    (Corollary~\ref{cor:strip-mixing}). The collapse lemma reduces the
    problem to a heat-bath interface chain, which a concentration-driven
    coupling mixes in two sweeps.
    Irreducibility on rectangles (Conjecture~\ref{conj:irred}) and mixing
    on two-dimensional grids remain open.
  \item \textbf{Machine verification (Section~\ref{sec:computation}).} All
    algebraic results were verified exhaustively on small graphs, including
    the torsor axioms on multigraph quotients, and we report an exhaustive
    irreducibility computation for a boundary-restricted chain on the
    $3 \times 3$ grid. We also report a computation showing that partition
    fibers are \emph{not} cosets of subgroups of $\Pick{G}$, evidence that
    the collapse route of (C3), rather than a group-theoretic treatment of
    balance, is the right division of labor.
\end{enumerate}

\done{The introduction is essentially complete; the related-work paragraphs
below may need expansion after the missing results stabilize.}

\subsection{Overview of the argument}
\label{subsec:story}

Sampling a balanced districting plan amounts to running a random walk
on the balanced $k$-forests of a graph. Two separate effects slow every
walk of this kind. Forests that induce the same district map differ in
their within-district trees, and a walk spends most of its steps
reshuffling those trees. Beyond that, the walk has to move between maps,
and the geometry of the balanced maps may itself be difficult to cross.

We treat the two effects separately. The first is removed exactly.
Within a fixed map the within-district trees correspond to the spanning
trees of the graph obtained by collapsing each district to a point, and
the sandpile group of the collapsed graph acts on them freely and
transitively (Sections~\ref{sec:torsor} and~\ref{sec:fiber}). The
hidden part of the state can therefore be redrawn exactly, in one step,
by Wilson's algorithm. The fiber collapse lemma of
Section~\ref{sec:collapse} shows that a chain which redraws the hidden
part exactly mixes precisely as fast as its projection onto maps, so
the entire problem reduces to the projected walk.

We then analyze the projected walk in a setting where it can be
understood completely, the $2 \times n$ strip. A balanced map of the
strip is described by the widths of the $k-1$ interfaces between
consecutive districts, and the target distribution favors narrow
interfaces, at a cost of roughly $\lambda = 2 + \sqrt3$ per additional
cell of width (Section~\ref{subsec:strips}). Redrawing one interface is
consequently a heavily biased draw, and two coupled copies of the walk
agree after each interface has been redrawn twice
(Section~\ref{subsec:coupling}). This yields the first Markov chain on
balanced partitions with a proven polynomial mixing time
(Corollary~\ref{cor:strip-mixing}). Square grids remain open, and
Section~\ref{sec:computation} records the estimates a two-dimensional
proof would require.

\subsection{Related work}

Torsor structures of the sandpile group on \emph{spanning trees} originate
with rotor-routing~\cite{holroyd2008chip} and the Bernardi
process~\cite{baker2018bernardi}; on plane graphs the two coincide and are
canonical~\cite{baker2018bernardi, chan2015rotor}, and rotor-routing is the
\emph{only} consistent torsor algorithm there~\cite{ganguly2023rotor}. Off
planarity both actions depend on the base data~\cite{shokrieh2021torsor}.
Our Theorem~\ref{thm:torsor} is a reduction to~\cite{holroyd2008chip}; its
value is the systematic transfer of the tree-level theory to rooted
$k$-forests and the resulting samplers. On the redistricting side, the
irreducibility state of the art is~\cite{akitaya2026bud} (exact balance,
$k = 2$, simple grid graphs, and triomino districts on rectangles); the
mixing state of the art for unconstrained forests
is~\cite{anari2021logconcave}; balanced sampling on grids by rejection
is~\cite{cannon2024balanced}; and the enumerative structure of balanced
trees of $2 \times n$ strips is~\cite{gallagher2025balanced}. The
decomposition method for Markov chains goes back
to~\cite{madras2002decomposition}; our Lemma~\ref{lem:collapse} is the
degenerate (and, for our purposes, optimal) case in which the
restriction chains are exact resamplers, so that the decomposition becomes an
identity rather than an inequality.

%═══════════════════════════════════════════════════════════════════════════════
\section{Preliminaries}
\label{sec:prelim}
%═══════════════════════════════════════════════════════════════════════════════

A handful of symbols carry the paper. In
Sections~\ref{sec:torsor}--\ref{sec:collapse}, $G$ is any connected
graph and $n = |V|$. In the strip sections the graph is the
$2 \times n$ grid, whose $2n$ cells form $k$ districts of $s = 2n/k$
cells each. Throughout, $t(H)$ is the number of spanning trees of $H$,
and $\lambda = 2 + \sqrt 3$ is the growth rate of $t$ on strips.
Everything else is introduced where it is first used.

\subsection{Graphs, forests, partitions}

All graphs are finite and connected; multigraphs (parallel edges) are allowed
and arise naturally as quotients; loops are deleted whenever they arise. For
$U \subseteq V$, $G[U]$ is the induced subgraph. A \emph{spanning
$k$-forest} is an acyclic edge set $F \subseteq E$ with exactly $k$
connected components whose vertex classes we denote $\xi(F) = \{\xi_1,
\ldots, \xi_k\}$; then $|F| = n - k$ and every $\xi_i$ induces a connected
subgraph, so $\xi(F)$ is a connected $k$-partition. Conversely the set of
$k$-forests inducing a fixed $\xi$ is $\prod_i \Tcal(G[\xi_i])$, where
$\Tcal(H)$ denotes the set of spanning trees of $H$; hence the uniform
distribution on $k$-forests projects to $\pi_{\mathrm{sp}}(\xi) \propto
\prod_i t(G[\xi_i])$ on connected $k$-partitions.

Given a connected $k$-partition $\xi$, the \emph{quotient multigraph}
$Q_\xi$ has one vertex per district and one edge per \emph{boundary edge} of
$G$ (an edge with endpoints in two different districts), with multiplicities
retained. A spanning tree $T \in \Tcal(G)$ is \emph{compatible} with $\xi$
if $T \cap E(G[\xi_i])$ is a spanning tree of $G[\xi_i]$ for every $i$; in
that case $T$ contains exactly $k-1$ boundary edges and these project to a
spanning tree of $Q_\xi$.

Fix a nonnegative vertex weight (population) $p : V \to \R_{\ge 0}$, with
$p \equiv 1$ unless stated otherwise, and $\bar{p} = p(V)/k$. A connected
$k$-partition is \emph{$\varepsilon$-balanced} if
$|p(\xi_i) - \bar{p}| \le \varepsilon \bar{p}$ for all $i$; ``exact
balance'' is $\varepsilon = 0$. Write $\Bcal_k^\varepsilon(G)$ for the set
of $\varepsilon$-balanced connected $k$-partitions and
$\Fcal_k^\varepsilon(G)$ for the set of $k$-forests $F$ with
$\xi(F) \in \Bcal_k^\varepsilon(G)$.

\subsection{Divisors, the Picard group, and reduced Laplacians}

$\operatorname{Div}(G) = \Z^V$ is the group of \emph{divisors} (chip
configurations); $\Delta_G$ is the graph Laplacian, acting on $\Z^V$ with
$\Delta_G(v,v) = \deg v$ and $\Delta_G(u,v) = -(\text{\# edges } uv)$.
Divisors $D \sim D'$ are \emph{linearly equivalent} if
$D - D' \in \im \Delta_G$. The \emph{Picard (sandpile, critical) group} of a
connected multigraph $H$ is
$\PicO{H} = \operatorname{Div}^0(H) / \im \Delta_H$, where
$\operatorname{Div}^0$ is the degree-zero subgroup; by the matrix-tree
theorem $|\PicO{H}| = t(H)$. For $W \subseteq V$ the matrix $\Delta^{(W)}$
is $\Delta_G$ with the rows and columns of $W$ deleted. The
\emph{Matrix-Forest theorem} (principal-minors form; see
\cite{chebotarev2006forest}) states that for $|W| = k$,
\[
\begin{gathered}
  \det \Delta^{(W)}
    \;=\; \#\bigl\{\text{spanning $k$-forests of $G$ in which} \\
    \text{each component contains exactly one vertex of } W\bigr\}.
\end{gathered}
\]
We call such forests \emph{rooted at $W$} and write $\Fcal_W(G)$ for
their set. We set
\[
  \Pick{G} \;=\; \Z^{n-k} \,/\, \im \Delta^{(S)}
\]
for a fixed root set $S$ of size $k$, so $|\Pick{G}| = |\Fcal_S(G)|$.

\subsection{Rotor-routing}

Let $H$ be a connected multigraph with a distinguished \emph{sink} $s$, and
fix a \emph{ribbon structure}: a cyclic ordering of the edge-ends at each
vertex. A \emph{rotor configuration} assigns to each non-sink vertex $v$ one
of its incident edge-ends $\rho(v)$. Adding a chip at $v$ and
\emph{rotor-routing} means: while the chip is at a non-sink vertex $u$,
advance $\rho(u)$ to the next incident edge-end in the cyclic order and move
the chip along that edge; stop when the chip reaches $s$. If the initial
rotors $\{(v, \rho(v))\}$ form a spanning tree oriented toward $s$, so do the
final rotors. The fundamental fact we use is:

\begin{fact}[{Holroyd--Levine--M\'esz\'aros--Peres--Propp--Wilson
\cite{holroyd2008chip}; see also \cite[Thm.~2.18]{ganguly2023rotor}}]
\label{fact:hlmppw}
For every connected multigraph $H$ with sink $s$ and every ribbon structure,
the rule ``add a chip at $v$ and rotor-route to the sink'' induces a
well-defined action of $\PicO{H}$ on $\Tcal(H)$ which is free and
transitive.
\end{fact}

\subsection{Markov chain conventions}

All chains are on finite state spaces. A kernel $P$ is reversible with
respect to $\pi$ if $\pi(x) P(x,y) = \pi(y) P(y,x)$; then $P$ is
self-adjoint on $L^2(\pi)$ with real spectrum
$1 = \lambda_1 \ge \lambda_2 \ge \cdots \ge \lambda_{|\Omega|} \ge -1$. The
\emph{spectral gap} is $\gamma(P) = 1 - \lambda_2(P)$, and
$t_{\mathrm{mix}}(P, \epsilon) = \min\{t : \max_{x}
\tv{P^t(x, \cdot) - \pi} \le \epsilon\}$.

\done{Section~\ref{sec:prelim} is complete and self-contained.}

%═══════════════════════════════════════════════════════════════════════════════
\section{The multi-sink sandpile torsor}
\label{sec:torsor}
%═══════════════════════════════════════════════════════════════════════════════

Throughout this section fix a root set $S = \{q_1, \ldots, q_k\} \subseteq
V$, $2 \le k < n$. Let $G_S$ be the multigraph obtained from $G$ by
identifying all vertices of $S$ to a single vertex $s^*$ and deleting the
loops created by edges with both endpoints in $S$. Edges of $G$ not inside
$S$ correspond bijectively to edges of $G_S$; we use this identification of
edge sets silently.

\begin{theorem}[Multi-sink torsor]
\label{thm:torsor}
Let $G$ be connected and $S$ as above.
\begin{enumerate}[(i)]
  \item $\Delta^{(S)}$ equals the reduced Laplacian of $G_S$ at $s^*$. In
    particular $\Pick{G} \cong \PicO{G_S}$, an isomorphism of explicitly
    presented groups (identical integer matrices, hence identical Smith
    normal forms), and $|\Pick{G}| = t(G_S)$.
  \item The map $F \mapsto F$ (on edge sets) is a bijection
    $\Fcal_S(G) \to \Tcal(G_S)$ between spanning $k$-forests of $G$ rooted
    at $S$ and spanning trees of $G_S$.
  \item For every ribbon structure on $G_S$, rotor-routing to the sink
    $s^*$ defines a free and transitive action of $\Pick{G}$ on
    $\Fcal_S(G)$.
\end{enumerate}
\end{theorem}

Informally, rooted $k$-forests of $G$ are the spanning trees of the
graph obtained by identifying the $k$ roots, and the sandpile group of
that graph permutes them freely and transitively. A uniform rooted
forest can therefore be drawn exactly, and no Markov chain is needed.

\begin{proof}
(i) The vertices of $G_S$ other than $s^*$ are exactly $V \setminus S$. For
$u, v \in V \setminus S$, identification changes neither the multiplicity of
edges between $u$ and $v$ nor $\deg u$: every edge from $u$ into $S$ becomes
an edge from $u$ to $s^*$, and no incident edge of $u$ is deleted (only
edges inside $S$ become loops). Hence the reduced Laplacian of $G_S$ at
$s^*$, the matrix indexed by $V \setminus S$ with diagonal $\deg_{G}(u)$
and off-diagonal $-(\text{\# edges } uv \text{ in } G)$, coincides entry
by entry with $\Delta^{(S)}$. The group isomorphism and the counting
consequence follow from the definitions and the matrix-tree theorem for
multigraphs.

(ii) Let $F \in \Fcal_S(G)$. No edge of $F$ joins two vertices of $S$: such
an edge would place two roots in one component. Hence $F$ is a well-defined
edge set of $G_S$. It has $n - k$ edges and $G_S$ has $n - k + 1$ vertices,
so it suffices to show the image is connected and acyclic. Connected:
identifying the $k$ roots glues the $k$ component trees of $F$ at the single
vertex $s^*$. Acyclic: a cycle in the image either avoids $s^*$, lifting to
a cycle of $F$ in $G$ (impossible), or passes through $s^*$, lifting to a
path in $F$ between two distinct roots (impossible: distinct roots lie in
distinct components of $F$).

Conversely let $T \in \Tcal(G_S)$ and let $F$ be the corresponding edge set
of $G$. $F$ is acyclic: a cycle of $F$ avoiding $S$ maps to a cycle of $T$;
a path of $F$ between two distinct roots maps to a cycle of $T$ through
$s^*$; a cycle of $F$ through one root maps to a cycle of $T$. Having
$n - k$ edges on $n$ vertices, $F$ is a spanning $k$-forest. Every
$v \in V \setminus S$ is connected in $F$ to some root: the $T$-path from
$v$ to $s^*$ lifts, edge by edge until its first edge incident to $s^*$, to
an $F$-path from $v$ to some $q_i$. So each of the $k$ components of $F$
contains at least one root, and by counting exactly one. The two maps are
mutually inverse because both are the identity on edge sets.

(iii) Combine (i), (ii), and Fact~\ref{fact:hlmppw} applied to $H = G_S$
with sink $s^*$.
\end{proof}

\done{Theorem~\ref{thm:torsor} is fully proved above. It was additionally
verified by machine (freeness, transitivity, commutativity, and the
chip-firing relations of the action) on the $2 \times 3$ grid with root sets
of sizes $2$ and $3$, the $3 \times 3$ grid, and $K_4$ with an adjacent root
pair, where $G_S$ has parallel edges; see Section~\ref{sec:computation}.}

\begin{remark}[Multi-sink Dhar burning]
\label{rem:dhar}
Dhar's burning algorithm on $G_S$ started at $s^*$ is precisely the
multi-source process on $G$ in which fire spreads simultaneously from all
$k$ roots; the edges that carry the fire form the rooted $k$-forest, and the
root whose fire claims a vertex defines a Voronoi-type assignment. All
recurrence tests and reduced-divisor computations for the $k$-sink sandpile
on $G$ are ordinary single-sink computations on $G_S$.
\end{remark}

\begin{remark}[Canonicity caveat]
\label{rem:canonicity}
$G_S$ is in general non-planar even when $G$ is planar, so the action of
Theorem~\ref{thm:torsor}(iii) depends on the ribbon structure and there is
no canonical choice: off planarity both the rotor and Bernardi torsors are
base-dependent~\cite{shokrieh2021torsor}, and on plane graphs rotor-routing
is the \emph{only} consistent torsor algorithm~\cite{ganguly2023rotor}. None
of the applications below require canonicity.
\end{remark}

\begin{corollary}[Exact sampling and navigation]
\label{cor:sampling}
Fix $S$.
\begin{enumerate}[(a)]
  \item Running Wilson's algorithm~\cite{wilson1996generating} on $G_S$
    rooted at $s^*$ returns a uniformly random element of $\Fcal_S(G)$ in
    expected time proportional to the mean hitting time of $G_S$ (in
    particular, expected polynomial time; $O(n \log n)$ on bounded-degree
    planar graphs such as grid quotients).
  \item Given the Smith normal form
    \[ U \Delta^{(S)} V = \operatorname{diag}(d_1, \ldots, d_{n-k}), \]
    a uniform element of $\Pick{G} \cong \prod_i \Z/d_i$ is sampled by
    choosing uniform coordinates, and its action on a stored reference
    forest is evaluated in polynomial time by decomposing a degree-zero
    representative into elementary chip additions and rotor-routing each
    chip to the sink.
\end{enumerate}
\end{corollary}

\begin{proof}
(a) is Wilson's theorem applied to $G_S$, using
Theorem~\ref{thm:torsor}(ii) to transport uniformity. (b) The SNF provides
the isomorphism and a representative divisor for each coordinate vector; a
degree-zero divisor $D$ on $G_S$ can be written $D = \sum_v D(v) \bigl((v) -
(s^*)\bigr)$ with $\sum_v |D(v)|$ bounded polynomially after replacing $D$
by its $s^*$-reduced representative (computable by Dhar's algorithm in
polynomial time), and each elementary action $(v) - (s^*)$ is one rotor walk
of at most $2|E|$ steps by \cite[Prop.~3.3]{ganguly2023rotor} when $v$ is
adjacent to $s^*$, and at most polynomially many steps in general
\cite{holroyd2008chip}.
\end{proof}

\statusnote{Corollary~\ref{cor:sampling}(b) is stated in the safe direction:
we sample a group element and \emph{apply} it. The inverse problem, that of computing from two forests the unique group
element carrying one to the other under the rotor action (the torsor discrete
logarithm), is not needed
anywhere in this paper. If we later want it for canonical-path
constructions, the clean route is to use any polynomial-time bijection
$\Tcal(G_S) \leftrightarrow \PicO{G_S}$ (e.g.\ Dhar's bijection through
reduced divisors) instead of inverting the rotor action; verify carefully
which statements need the action itself and which need only \emph{some}
bijection.}

%═══════════════════════════════════════════════════════════════════════════════
\section{Fibers and exact fiber resampling}
\label{sec:fiber}
%═══════════════════════════════════════════════════════════════════════════════

Fix a connected $k$-partition $\xi$. Two state spaces occur in the
literature: \emph{forest-valued} chains store $F \in \Fcal_k(G)$, and
\emph{tree-valued} chains store a compatible spanning tree (possibly with
marked boundary edges). We treat both; the tree case is where the quotient
group enters.

\begin{theorem}[Fiber decomposition]
\label{thm:fiber-decomp}
For every connected $k$-partition $\xi$, the map
\[
  T \;\longmapsto\;
  \Bigl( \bigl(T \cap E(G[\xi_i])\bigr)_{i=1}^{k},\;
         T \cap \partial\xi \Bigr)
\]
is a bijection between the compatible spanning trees of $G$ and
$\prod_{i=1}^k \Tcal(G[\xi_i]) \times \Tcal(Q_\xi)$. Consequently
\[
\begin{gathered}
  \tau(\xi) := \#\{\text{compatible spanning trees}\}
   = \Bigl(\prod_{i=1}^{k} t(G[\xi_i])\Bigr) \cdot t(Q_\xi), \\[4pt]
  \#\{\text{$k$-forests with } \xi(F) = \xi\} = \prod_{i=1}^{k} t(G[\xi_i]).
\end{gathered}
\]
\end{theorem}

\begin{proof}
If $T$ is compatible, each $T \cap E(G[\xi_i])$ is by definition a spanning
tree of $G[\xi_i]$, accounting for $\sum_i (|\xi_i| - 1) = n - k$ edges; the
remaining $k - 1$ edges of $T$ are boundary edges, and they connect the $k$
districts without cycles (contracting each district maps $T$ onto a
connected, $(k-1)$-edge subgraph of $Q_\xi$), i.e.\ they form a spanning
tree of the multigraph $Q_\xi$. Conversely, the union of spanning trees
$T_i$ of the districts and a spanning tree $\sigma$ of $Q_\xi$ (a choice of
$k-1$ boundary edges of $G$ connecting all districts acyclically) has
$n - 1$ edges, spans $V$, and is acyclic: a cycle would either live inside a
district (impossible, $T_i$ is a tree) or project to a nontrivial closed
walk in the tree $\sigma$ (impossible). The two constructions are mutually
inverse. The forest statement is the case in which no boundary edges are
selected.
\end{proof}

\done{Theorem~\ref{thm:fiber-decomp} is fully proved and was verified by
brute-force enumeration on five partitions of the $2\times 3$ and
$3 \times 3$ grids and on a $6 \times 6$ instance
(Section~\ref{sec:computation}).}

\begin{proposition}[Exact fiber resampler]
\label{prop:resampler}
Fix $\xi$. Each of the following runs in expected polynomial time and
returns an exactly uniform sample:
\begin{enumerate}[(a)]
  \item a uniform $k$-forest with partition $\xi$: run Wilson's algorithm
    independently in each $G[\xi_i]$;
  \item a uniform compatible spanning tree: additionally sample a uniform
    spanning tree of the multigraph $Q_\xi$, either by Wilson's algorithm
    on $Q_\xi$, or by sampling a uniform $g \in \PicO{Q_\xi}$ via the Smith
    normal form and letting $g$ act on a reference tree of $Q_\xi$ through
    rotor-routing (Fact~\ref{fact:hlmppw}), and take the union as in
    Theorem~\ref{thm:fiber-decomp}.
\end{enumerate}
\end{proposition}

\begin{proof}
Immediate from Theorem~\ref{thm:fiber-decomp}, Wilson's theorem, and, for
the torsor route in (b), the fact that if the action of a finite group on a
finite set is free and transitive then the orbit of any fixed point under a
uniform group element is uniform.
\end{proof}

\begin{remark}
The torsor route in (b) is stated because the group $\PicO{Q_\xi}$, not
any particular sampler, is the structural object: it indexes the
\emph{boundary configurations} of $\xi$ and quotients away the within-fiber
redundancy that tree-valued chains otherwise spend most of their steps
shuffling~\cite{mew}. In implementations Wilson on $Q_\xi$ is simpler and
equally exact.
\end{remark}

%═══════════════════════════════════════════════════════════════════════════════
\section{The fiber collapse lemma}
\label{sec:collapse}
%═══════════════════════════════════════════════════════════════════════════════

We now formalize the removal of (S1). The lemma is stated
abstractly; the intended instance is $\Omega = \Fcal_k^\varepsilon(G)$ (or
the compatible-tree space), fibered over
$\Xi = \Bcal_k^\varepsilon(G)$ by $x \mapsto \xi(x)$, with $P_F$ the exact
resampler of Proposition~\ref{prop:resampler}.

\paragraph{Setting.} Let $\Omega$ be a finite set, $\zeta : \Omega \to \Xi$
a surjection onto a finite index set (\emph{fibration}), and $\pi$ a
probability distribution on $\Omega$ with full support. Write
$\bar\pi = \zeta_* \pi$ for the pushforward and
$\nu_\xi = \pi(\,\cdot \mid \zeta = \xi)$ for the conditionals, so that
$\pi = \sum_{\xi} \bar\pi(\xi)\, \nu_\xi$ and the $\nu_\xi$ have pairwise
disjoint supports. Define the \emph{fiber resampling kernel}
\[
  P_F(x, y) = \nu_{\zeta(x)}(y),
\]
and let $P_M$ be any transition kernel on $\Omega$ reversible with respect
to $\pi$. Define the \emph{sandwich kernel} $P = P_F P_M P_F$ (apply $P_F$,
then $P_M$, then $P_F$) and the \emph{projected kernel} on $\Xi$
\[
  \bar{P}(\xi, \xi') \;=\; \sum_{x \in \zeta^{-1}(\xi)} \nu_\xi(x)
    \sum_{y \in \zeta^{-1}(\xi')} P_M(x, y).
\]

\begin{lemma}[Fiber collapse]
\label{lem:collapse}
With the notation above:
\begin{enumerate}[(a)]
  \item $P_F$ is reversible with respect to $\pi$, idempotent
    ($P_F^2 = P_F$), and as an operator on $L^2(\pi)$ equals the orthogonal
    projection $\Pi f = \E_\pi[f \mid \zeta]$ onto the subspace of
    $\zeta$-measurable functions. Consequently $P = \Pi M \Pi$ (with $M$
    the operator of $P_M$) is reversible with respect to $\pi$.
  \item $\bar P$ is reversible with respect to $\bar\pi$, and for the
    $P$-chain $(x_t)_{t \ge 0}$ started at any $x_0$, the projected process
    $\xi_t = \zeta(x_t)$ is a Markov chain with kernel $\bar P$ started at
    $\zeta(x_0)$.
  \item For every $t \ge 1$ and every $x_0$,
    \[
      \tv{\,\mathrm{Law}(x_t) - \pi\,}
        \;=\; \tv{\,\mathrm{Law}(\xi_t) - \bar\pi\,}
        \;=\; \tv{\,\bar P^{\,t}(\zeta(x_0), \cdot) - \bar\pi\,}.
    \]
    In particular $t_{\mathrm{mix}}(P, \epsilon) =
    t_{\mathrm{mix}}(\bar P, \epsilon)$ for every $\epsilon$.
  \item $\operatorname{Spec}(P) = \operatorname{Spec}(\bar P) \cup \{0\}$
    (as multisets, the eigenvalue $0$ appearing with multiplicity
    $|\Omega| - |\Xi|$). In particular
    $\lambda_2(P) = \max\{\lambda_2(\bar P), 0\}$ and
    $\lambda_{\min}(P) = \min\{\lambda_{\min}(\bar P), 0\}$.
\end{enumerate}
\end{lemma}

Informally, once the hidden coordinate is redrawn exactly at every
step it no longer matters; the chain and its projection onto partitions
are equally far from stationarity at every time.

\begin{proof}
(a) For $x, y$ in the same fiber $\zeta^{-1}(\xi)$,
$\pi(x) P_F(x,y) = \bar\pi(\xi) \nu_\xi(x) \nu_\xi(y)$ is symmetric in
$(x,y)$; for $x, y$ in different fibers both sides vanish. Idempotence is
the fact that resampling twice from the same conditional is the same as
resampling once ($\zeta$ is unchanged by $P_F$). For the operator
statement, $(P_F f)(x) = \sum_y \nu_{\zeta(x)}(y) f(y) =
\E_\pi[f \mid \zeta = \zeta(x)]$, which is precisely the conditional
expectation; conditional expectation is the orthogonal projection onto
$L^2(\bar\pi) \subseteq L^2(\pi)$, the subspace of $\zeta$-measurable
functions. Then $P^* = P_F^* P_M^* P_F^* = P_F P_M P_F = P$ since $P_F$,
$P_M$ are self-adjoint.

(b) Reversibility of $\bar P$: for $g, h : \Xi \to \R$,
$\langle \bar P g, h \rangle_{\bar\pi} = \langle M (g \circ \zeta),
h \circ \zeta \rangle_{\pi}$ is symmetric because $M$ is self-adjoint.
Markov property: one step of $P$ from $x_t$ consists of $y \sim
\nu_{\zeta(x_t)}$, then $z \sim P_M(y, \cdot)$, then $x_{t+1} \sim
\nu_{\zeta(z)}$. The law of $\zeta(x_{t+1}) = \zeta(z)$ depends on the past
only through $\zeta(x_t)$ (the draw $y$ erases the position within the
fiber), and equals $\bar P(\zeta(x_t), \cdot)$ by definition of $\bar P$.

(c) After the terminal $P_F$ in step $t$, the conditional law of $x_t$
given $(\xi_0, \ldots, \xi_t)$ is exactly $\nu_{\xi_t}$. Hence
$\mathrm{Law}(x_t) = \sum_\xi \Pr[\xi_t = \xi]\, \nu_\xi$, while $\pi =
\sum_\xi \bar\pi(\xi) \nu_\xi$. Because the $\nu_\xi$ have disjoint
supports,
\[
  \tv{\mathrm{Law}(x_t) - \pi}
   = \tfrac12 \sum_{\xi} \sum_{x \in \zeta^{-1}(\xi)}
     \bigl| \Pr[\xi_t = \xi] - \bar\pi(\xi) \bigr|\, \nu_\xi(x)
   = \tfrac12 \sum_{\xi} \bigl| \Pr[\xi_t = \xi] - \bar\pi(\xi) \bigr|,
\]
which is the projected total variation distance; the second equality in the
statement is (b). Taking maxima over starting states (every $\xi_0 \in \Xi$
is realized by some $x_0$) gives the equality of mixing times.

(d) $L^2(\pi) = \operatorname{ran} \Pi \oplus \ker \Pi$. On
$\operatorname{ran} \Pi \cong L^2(\bar\pi)$ we have, for $f = g \circ
\zeta$: $P f = \Pi M \Pi f = \Pi M f = (\bar P g) \circ \zeta$, so
$P|_{\operatorname{ran}\Pi}$ is conjugate to $\bar P$. On $\ker \Pi$:
$P f = \Pi M \Pi f = \Pi M 0 = 0$. The spectral statement follows.
\end{proof}

\done{Lemma~\ref{lem:collapse} is fully proved above and is elementary; its
value is that Proposition~\ref{prop:resampler} makes $P_F$ \emph{exactly}
implementable for balanced forests, which is what turns the classical
decomposition inequality of~\cite{madras2002decomposition} into an identity.}

\begin{theorem}[Conditional mixing theorem]
\label{thm:conditional}
Let $\Omega = \Fcal_k^\varepsilon(G)$ with $\zeta = \xi(\cdot)$, let $\pi$
be any distribution of the product form $\pi(F) = \bar\pi(\xi(F))\,
\nu_{\xi(F)}(F)$ with $\nu_\xi$ uniform on the fiber (e.g.\ $\pi$ uniform on
$\Fcal_k^\varepsilon(G)$, for which $\bar\pi = \pi_{\mathrm{sp}}$
conditioned on balance), and let $P_M$ be any $\pi$-reversible
partition-changing kernel with each step implementable in polynomial time.
Then the sandwich chain $P = P_F P_M P_F$ is implementable in expected
polynomial time per step, reversible with respect to $\pi$, and satisfies
\[
  t_{\mathrm{mix}}(P, \epsilon) = t_{\mathrm{mix}}(\bar P, \epsilon)
  \quad \text{for all } \epsilon .
\]
In particular, if the projected chain $\bar P$ on
$\Bcal_k^\varepsilon(G)$ is irreducible with spectral gap
$\gamma(\bar P) \ge 1/\mathrm{poly}(n)$, then $P$ mixes in polynomial time
on the balanced forest space. The same statements hold verbatim for the
compatible-tree state space using
Proposition~\ref{prop:resampler}(b).
\end{theorem}

\begin{proof}
Combine Proposition~\ref{prop:resampler} (implementability of $P_F$),
Lemma~\ref{lem:collapse}(a)--(d), and the standard bound
$t_{\mathrm{mix}}(\bar P, \epsilon) \le
\gamma(\bar P)^{-1} \log \bigl( (\min_\xi \bar\pi(\xi))^{-1}
\epsilon^{-1} \bigr)$ for reversible chains (applied after the usual lazy
modification if $\lambda_{\min}(\bar P) < 0$; laziness changes the gap by a
factor $2$ and survives the sandwich construction unchanged).
\end{proof}

\done{Theorem~\ref{thm:conditional} is proved in full, \emph{conditionally
on a gap bound for the projected chain}. It is the precise interface between
the algebraic half of the paper (done) and the geometric half (open,
Section~\ref{sec:targets}).}

\begin{remark}[Why not treat balance algebraically?]
One might hope the balanced set is itself an orbit or coset for a subgroup
of $\Pick{G}$, so that Section~\ref{sec:torsor} alone would finish the
problem. It is not: machine verification on the $2 \times 3$ grid shows that
already the fiber of a single partition is in general \emph{not} a coset of
any subgroup of $\Pick{G}$, nor is the union of the balanced fibers
(Section~\ref{sec:computation}). The group sees the forest structure
exactly and the balance structure not at all; Lemma~\ref{lem:collapse} is
the correct division of labor.
\end{remark}

%═══════════════════════════════════════════════════════════════════════════════
\section{The projected chain and the strip theorem}
\label{sec:targets}
%═══════════════════════════════════════════════════════════════════════════════

By Theorem~\ref{thm:conditional} everything now depends on the
projected kernel $\bar P$ on $\Bcal_k^\varepsilon(G)$. This section
fixes a concrete $P_M$, proves the main theorem on $2 \times n$ strips,
and states the rectangular case as a conjecture with a proof plan.

\subsection{The partition-changing kernel}
\label{subsec:pm}

We take $P_M$ to be the \emph{window resampling} kernel. From a forest
$F$ with partition $\xi$, choose an edge $\{i, j\}$ of $Q_\xi$ uniformly
at random and replace the forest on $H = G[\xi_i \cup \xi_j]$ by an exact
draw from its conditional distribution given the rest of the state; if
the window is disconnected, hold. Under the uniform distribution on
balanced forests the conditional factors. First choose a balanced split
$\{A, B\}$ of $H$ with probability proportional to $t(G[A])\, t(G[B])$,
then draw uniform spanning trees of $A$ and of $B$ by Wilson's
algorithm. On strips the split distribution is explicitly computable,
since Proposition~\ref{prop:windows} lists the $O(s)$ balanced splits
and their weights are products of ladder counts, available from the
three-term recurrence. $P_M$ is a Gibbs block update of the window
forest given the rest, hence reversible with respect to $\pi$.

\statusnote{This kernel replaced an earlier tree-based definition
(uniform spanning tree of the window, uniform balanced split edge, MH
correction). The pre-submission audit found two problems with the
tree-based route: the exact MH correction of the uniform-tree proposal
is not computable in closed form, and every implementable variant
conditions the tree through the current split, which under exact balance
freezes the projected chain outright
(Remark~\ref{rem:tree-kernels}; machine-verified). The heat-bath kernel
has neither problem on strips, and the collapse lemma is indifferent to
the choice.}

\subsection{Irreducibility (Conjecture A)}

\begin{conjecture}[Irreducibility, exact balance, all $k$]
\label{conj:irred}
Let $G$ be an $a \times b$ rectangular grid graph with $k \mid ab$. The
projected window-resampling chain $\bar P$ is irreducible on
$\Bcal_k^0(G)$, with diameter $O(ab)$ moves.
\end{conjecture}

\paragraph{Status and plan.} For $k = 2$ this follows for the BUD chain
from \cite[Thm.~B.2]{akitaya2026bud} on all simple grid graphs (any
prescribed part sizes, diameter $O(n)$); a transfer argument from BUD moves
to window moves is needed and appears routine, since every BUD move
between exactly balanced states is realized by one window resample of the
two districts involved (via \cite[Lemma~A.3]{akitaya2026bud}, retreeing is free
inside our framework: the fiber coordinate is irrelevant after collapse).
For $k \ge 3$ the natural strategy, proposed in~\cite{akitaya2026bud} and
currently open, is recursion: prove the $k = 2$ statement not just for
simple grid graphs but for the district-pair subgraphs
$G[\xi_i \cup \xi_j]$ that arise along the recursion; these are connected
but can have holes and cut vertices, where counterexamples are known to
exist in general \cite[Lemma~B.14]{akitaya2026bud}. The plan:
\begin{enumerate}[1.]
  \item Prove Conjecture~\ref{conj:irred} for $k = 3$ on rectangles
    directly, by adapting the canonicalization argument of
    \cite[App.~B]{akitaya2026bud}: canonicalize to bands of columns, using
    window moves of \emph{adjacent} bands only (the union of two adjacent
    canonical bands is simply connected, avoiding the holes obstruction).
  \item Identify the class of intermediate shapes closed under the
    recursion (candidate: ``column-convex'' district pairs) and prove the
    $k = 2$ statement for that class.
  \item General $k$ by induction, always merging a canonical extreme band
    with its neighbor.
\end{enumerate}

\missing{All three steps. Step 1 is a finite, technical extension of a
published argument and is the right first target ($\sim$weeks, not months).
Note the competition risk: the authors of~\cite{akitaya2026bud} state the
$k > 2$ extension as their own ongoing work. Our differentiator is that we
only need irreducibility for \emph{merge--split on rectangles}, not for BUD
on arbitrary grid subgraphs; check whether restricting the ambition this way
shortens their argument materially.}

Corollary~\ref{cor:strip-irred} proves the strip analogue of
Conjecture~\ref{conj:irred} outright; the triomino irreducibility of
\cite[Thm.~B.13]{akitaya2026bud} on rectangles and the $3 \times 3$
computation of Section~\ref{sec:computation} support the rectangular
case.

\statusnote{Empirical support: on the $3 \times 3$ grid with $k = 3$ and
exact balance, the \emph{partition-level} communication classes of a
boundary-restricted chain cover all $10$ balanced partitions
(Section~\ref{sec:computation}); and \cite[Thm.~B.13]{akitaya2026bud} proves
the triomino case ($k = ab/3$) on all rectangles, which includes
$3 \times 3$, $k = 3$.}

\subsection{The main theorem for $2 \times n$ strips}
\label{subsec:strips}

\begin{theorem}[Strip gap]
\label{conj:strip}
Fix $k \ge 3$ with $k \mid 2n$, exact balance, and write $s = 2n/k$. The
projected window-resampling chain $\bar P$ on $\Bcal_k^0(2 \times n)$
satisfies, with absolute constants $c, C_0 > 0$,
\[
  \gamma(\bar P) \;\ge\; \frac{c}{k \log(k+1)}
  \qquad \text{for } s \ge C_0 \log(k+1),
\]
and $\gamma(\bar P) > 0$ for the finitely many smaller $s$ (by
Corollary~\ref{cor:strip-irred}). The proof is in
Section~\ref{subsec:coupling}.
\end{theorem}

\begin{remark}[Why $k \ge 3$]
\label{rem:k2-trivial}
For $k = 2$ the only window is the whole strip, so a single update
redraws the entire partition from its exact conditional distribution;
the chain mixes in one step and yields no interesting theorem. The
substantive question, where updates are genuinely local, begins at
$k = 3$.
\end{remark}

\begin{corollary}[Strip mixing, unconditional]
\label{cor:strip-mixing}
Fix $k \ge 2$ with $k \mid 2n$ and $s = 2n/k$. Take $P_M$ to be the
window resampling kernel of Section~\ref{subsec:pm} and
$P = P_F P_M P_F$ the fiber-Gibbs balanced chain on
$\Fcal_k^0(2 \times n)$. Then
\[
  t_{\mathrm{mix}}(P, 1/4) \;\le\; 2(k-1)\bigl(\lceil \log(k+1)
  \rceil + 4\bigr)
  \qquad \text{for } s \ge C_0 \log(k+1),
\]
where $C_0$ is the constant of Lemma~\ref{lem:coupling}, and
$t_{\mathrm{mix}}(P, 1/4) \le C_k$ for the finitely many remaining $n$.
Each step runs in expected time $O(n \log n)$ measured in random
walk moves, plus arithmetic on $O(n)$-digit integers, so for every
fixed $k$ the sampler produces an approximately uniform balanced
$k$-forest in near-linear expected total time. To our knowledge this is
the first polynomial mixing-time bound for a Markov chain whose states
are
balanced connected partitions (lifted to forests) of a nontrivial graph
family. The proof is at the end of Section~\ref{subsec:coupling}.
\end{corollary}

\paragraph{The balanced partition space of a strip.} We settle the $k = 2$
structure completely; it is the template for the general interface coding.
Index the vertices of the $2 \times n$ grid by $\{0,1\} \times \{0, \ldots,
n-1\}$ (row, column); \emph{column} $c$ is the pair $\{(0,c), (1,c)\}$. For
a $2$-partition $\{A, B\}$, call a column \emph{pure-$A$}
(resp.\ \emph{pure-$B$}) if both its vertices lie in $A$ (resp.\ $B$), and
\emph{mixed} otherwise; the \emph{orientation} of a mixed column records
which part holds its top vertex.

\begin{proposition}[Classification of balanced $2$-partitions of the strip]
\label{prop:strip-classify}
Let $n \ge 2$. The balanced connected $2$-partitions of the $2 \times n$
grid are exactly the following, and there are exactly $n$ of them:
\begin{enumerate}[(a)]
  \item the horizontal split $\{$top row, bottom row$\}$;
  \item for $n$ even, the vertical split after column $n/2 - 1$;
  \item for each $a \in \{1, \ldots, \lceil n/2 \rceil - 1\}$ and each
    $\sigma \in \{\mathrm{top}, \mathrm{bottom}\}$, the \emph{staircase}
    $S_{a,\sigma} = \{P_{a,\sigma},\, V \setminus P_{a,\sigma}\}$, where
    $P_{a,\sigma}$ consists of the $a$ leftmost (full) columns together
    with the $\sigma$-row vertices of the middle $n - 2a$ columns.
\end{enumerate}
\end{proposition}

\begin{proof}
Each listed partition is balanced ($|P_{a,\sigma}| = 2a + (n-2a) = n$) and
has connected parts: the middle vertices of $P_{a,\sigma}$ form a
horizontal path attached to its full block, and the complement is the
left--right, top--bottom mirror image. The left--right mirror of
$S_{a,\mathrm{top}}$ coincides with $S_{a,\mathrm{bottom}}$ \emph{as an
unordered partition}, so each unordered partition appears exactly once in
the list, and the listed partitions are pairwise distinct. It remains to
show there are no others.

Let $\{A, B\}$ be a balanced connected $2$-partition, with $a', b', m$ the
numbers of pure-$A$, pure-$B$, and mixed columns. Every mixed column
contributes one vertex to each part, so $n = |A| = 2a' + m$ and
$n = |B| = 2b' + m$; hence $a' = b'$.

\emph{Step 1: no pure columns $\Rightarrow$ horizontal split.} If all
columns are mixed and two adjacent columns had different orientations,
$A$ would meet both rows; a connected set meeting both rows must contain a
vertical edge, i.e.\ a pure column, a contradiction. So all orientations
agree and $A$ is a full row.

\emph{Step 2: the pure-$A$ columns are consecutive; likewise pure-$B$.}
Suppose pure-$A$ columns $c < c'$ have a non-pure-$A$ column strictly
between them. Any path between vertices on opposite sides of column $c$
passes through column $c$, which lies entirely in $A$; hence the connected
set $B$ has no vertices on both sides of $c$, nor of $c'$. Since some
column strictly between $c$ and $c'$ contains a $B$-vertex, $B$ lies
entirely in the open ``pocket'' between $c$ and $c'$. Because $A$ is
connected, some $A$-path crosses the pocket from column $c$ to column
$c'$; inside the pocket it cannot change rows (a row change uses a
vertical edge with both endpoints in $A$, i.e.\ a pure-$A$ column, and the
pocket contains none), so one full row of the pocket lies in $A$. Every
pocket column, being non-pure-$A$, then has exactly its other-row vertex
in $B$, whence $|B| = (\text{pocket width}) \le n - 2 < n$, contradicting
balance.

\emph{Step 3: mixed columns lie between the two pure blocks, which sit at
the two ends.} If $a' = 0$ we are in Step 1, so let $a' = b' \ge 1$; by
Step 2 there is a single pure-$A$ block and a single pure-$B$ block. A
mixed column with both blocks on the same side of it has a $B$-vertex
(or, in the mirrored case, an $A$-vertex) that must reach its block by a
path crossing the other part's pure block, which blocks both rows; this
is impossible. Hence all mixed columns lie between the blocks, the pure-$A$
block starts at column $0$, and the pure-$B$ block ends at column $n-1$
(or the mirror image, which yields the same set of unordered partitions).

\emph{Step 4: the mixed run has constant orientation.} Let the mixed run
occupy columns $\ell, \ldots, r$ and suppose a flip occurs; let $d$ be the
leftmost mixed column whose orientation differs from that of column
$\ell$. Up to reflecting the rows, $A$ holds the tops of columns $\ell,
\ldots, d-1$ and the bottom of column $d$. Let $\{(1,d), \ldots, (1,e)\}$
be the maximal bottom-row run of $A$-vertices starting at $(1,d)$, with
$e \le r$. Its left neighbor $(1, d-1)$ is in $B$ (column $d-1$ has $A$
on top); its right neighbor $(1, e+1)$ is in $B$ (either a mixed column
with $A$ on top, by maximality, or the pure-$B$ block); and every vertical
neighbor $(0, j)$, $d \le j \le e$, is in $B$. So this run is a connected
component of $A$ not containing the pure-$A$ block, contradicting
connectivity.

Steps 1--4 leave exactly the pattern
$[\text{pure-}A]^{a} [\text{mixed, orientation } \sigma]^{\,n-2a}
[\text{pure-}B]^{a}$: this is $S_{a,\sigma}$ when $n - 2a \ge 1$, the
vertical split when $n = 2a$, and the horizontal split when $a = 0$.
Counting: $1 + 2(\lceil n/2 \rceil - 1) + \mathbf{1}_{\{n\ \mathrm{even}\}}
= n$ in both parities.
\end{proof}

\done{Proposition~\ref{prop:strip-classify} is proved in full and matches
the exhaustive machine counts for all $n \le 12$
(Section~\ref{sec:computation}, item 6). It also corrects an imprecision in
an earlier version of this section: the balanced states are staircases, not
column cuts; the vertical cut exists only for even $n$.}

\begin{proposition}[Stationary law on the strip, $k = 2$]
\label{prop:strip-law}
Let $t_a$ denote the number of spanning trees of the $2 \times a$ grid:
$t_1 = 1$, $t_2 = 4$, $t_a = 4t_{a-1} - t_{a-2}$. Under the uniform
distribution on balanced $2$-forests of the $2 \times n$ grid (equivalently,
$\pi_{\mathrm{sp}}$ conditioned on exact balance), the projected law on the
$n$ partitions of Proposition~\ref{prop:strip-classify} is
\[
  \bar\pi(\mathrm{horizontal}) \propto 1, \qquad
  \bar\pi(S_{a,\sigma}) \propto t_a^{\,2}, \qquad
  \bar\pi(\mathrm{vertical}) \propto t_{n/2}^{\,2}.
\]
In particular $\bar\pi$ is log-linear in the staircase depth $a$ with rate
$2\log(2 + \sqrt{3})$ and concentrates on the deepest cuts.
\end{proposition}

\begin{proof}
By Theorem~\ref{thm:fiber-decomp} the fiber of $\xi$ contains $\prod_i
t(G[\xi_i])$ forests. The parts of the horizontal split are paths, each
with one spanning tree. The induced graph of $P_{a,\sigma}$ is the
$2 \times a$ grid with a path of $n - 2a$ vertices attached by a single
edge at a corner; the attached path consists of bridges, present in every
spanning tree, so $t(G[P_{a,\sigma}]) = t_a$, and the complementary part is
isomorphic to its mirror image, giving $t_a^2$. The recurrence and growth
constant $2 + \sqrt 3$ (largest root of $x^2 = 4x - 1$) are classical for
strip graphs.
\end{proof}

\done{Proposition~\ref{prop:strip-law} is proved in full. Together the two
propositions complete Theorem~\ref{conj:strip} step 1 for $k = 2$ and
supply the tree-count calculus (district count $= t_{\mathrm{width}}$, since
staircase fringes attach as pendant paths at corners) used for the
conditionals in step 2.}

\begin{lemma}[Split counts and constant split probability on windows]
\label{lem:split}
Let $w \ge 2$, let $T$ be a uniform spanning tree of the $2 \times w$ grid,
and let $\xi = \{A, B\}$ range over the balanced connected $2$-partitions of
the window (Proposition~\ref{prop:strip-classify}).
\begin{enumerate}[(a)]
  \item The spanning trees with $|T \cap \partial\xi| = 1$ (equivalently,
    those containing an edge whose removal yields exactly $\xi$, necessarily
    unique) number exactly $t(G[A]) \cdot t(G[B]) \cdot |\partial\xi|$.
  \item For $w$ even,
    \[
      \Pr\bigl[T \text{ admits a balanced split}\bigr]
        \;\ge\; \frac{2\, t_{w/2}^2}{t_w}
        \;=\; \frac{1}{\sqrt 3}\cdot
          \frac{\lambda^w + \lambda^{-w} - 2}{\lambda^w - \lambda^{-w}}
        \;\ge\; \frac12,
      \qquad \lambda = 2 + \sqrt 3,
    \]
    with equality at $w = 2$ and limit $1/\sqrt 3$ as $w \to \infty$.
  \item A tree admits at most one balanced split edge per partition, hence
    at most $w$ balanced split edges in total. Consequently the
    tree proposal (a uniform spanning tree of the window followed
    by a uniform choice among its balanced split edges) proposes each
    balanced $\xi$ with probability at least
    $t(G[A])\, t(G[B])\, |\partial\xi| / (w \, t_w)$, which is within a
    factor $\mathrm{poly}(w)$ of the exact conditional
    $\bar\pi(\xi) \propto t(G[A])\, t(G[B])$: the move is a
    $\mathrm{poly}(w)$-accurate independence sampler for the interface
    conditional.
\end{enumerate}
\end{lemma}

\begin{proof}
(a) If $T \cap \partial\xi = \{e\}$, then $T \setminus \{e\}$ is an acyclic
subgraph of $G[A] \sqcup G[B]$ with $2w - 2$ edges, i.e.\ the disjoint union
of a spanning tree of $G[A]$ and one of $G[B]$; conversely every triple
$(T_A, T_B, e)$ with $e \in \partial\xi$ assembles to such a spanning tree,
and the two maps are mutually inverse. Removing an edge $e'$ of a spanning
tree produces the partition into the two components of $T \setminus \{e'\}$;
this equals $\xi$ if and only if $T \cap \partial\xi = \{e'\}$. So the
counted trees are exactly those admitting a split at $\xi$, with a unique
split edge.

(b) Apply (a) to the vertical cut: $G[A] \cong G[B] \cong 2 \times (w/2)$
grid and $|\partial\xi| = 2$, giving the lower bound $2 t_{w/2}^2 / t_w$.
With $t_a = (\lambda^a - \lambda^{-a})/(2\sqrt3)$ and
$\lambda \lambda^{-1} = 1$, the middle expression follows by expanding the
square. The inequality $\ge 1/2$ is equivalent to
$(1 - \tfrac{\sqrt3}{2})\lambda^w + (1 + \tfrac{\sqrt3}{2})\lambda^{-w}
\ge 2$, which holds with equality at $w = 2$ and whose left side is
increasing in $w \ge 2$ (its derivative has the sign of
$(1 - \tfrac{\sqrt3}{2})\lambda^w - (1 + \tfrac{\sqrt3}{2})\lambda^{-w}$,
positive already at $w = 2$).

(c) Uniqueness per partition is (a); Proposition~\ref{prop:strip-classify}
bounds the number of partitions by $w$. If a tree has $s \le w$ balanced
split edges, each is chosen with probability $1/s \ge 1/w$; combining with
(a) gives the proposal bound. Since $2 \le |\partial\xi| \le w + 2$, the
proposal law and $\bar\pi$ differ by factors in
$[\,2/(w(w+2)),\, 1\,] \cdot \mathrm{const}$.
\end{proof}

\done{Lemma~\ref{lem:split} is proved in full and machine-verified: the
counting identity of (a) holds exactly for every balanced partition at
$w = 4, 6, 8$ (exhaustive enumeration over $56$, $780$, $10{,}864$ trees);
the splittability probabilities are $0.786, 0.764, 0.763$ against the
bound $0.571, 0.577, 0.577$; and the measured proposal/target ratios lie
in $[0.86, 3.5]$. This retires the estimate previously flagged as the
technical heart of Theorem~\ref{conj:strip}: what remains of step 2 is
to extend the lemma from rectangular windows to the merged windows of
\emph{interior} district pairs, whose shape is a ladder with a pendant
staircase fringe at each end; the tree-count calculus is unchanged, as
pendant fringes contribute factor $1$.}

\paragraph{The general classification.} For $k \ge 3$ the same column
analysis settles the full structure. Write $s = 2n/k$, and set
$\hat t_0 = 1$ and $\hat t_a = t_a$ for $a \ge 1$ (the ladder counts of
Proposition~\ref{prop:strip-law}; a district with no full column is a
path and has one spanning tree).

\begin{proposition}[Classification, general $k$]
\label{prop:strip-classify-k}
Let $k \ge 2$ and $k \mid 2n$. The balanced connected $k$-partitions of
the $2 \times n$ grid are exactly the configurations produced as follows.
Choose integers $a_1, \ldots, a_k \ge 0$, set $\mu_0 = \mu_k = 0$ and
$\mu_i = s - 2a_i - \mu_{i-1}$ for $1 \le i \le k-1$, and require
$\mu_i \ge 0$ for all $i$ together with $2a_k = s - \mu_{k-1}$. The
column word, read left to right, is
\[
  P_1^{\,a_1}\, R_1^{\,\mu_1}\, P_2^{\,a_2}\, R_2^{\,\mu_2} \cdots
  R_{k-1}^{\,\mu_{k-1}}\, P_k^{\,a_k},
\]
where $P_i$ denotes a full column of district $i$ and each nonempty run
$R_j^{\,\mu_j}$ consists of mixed columns carrying districts $j$ and
$j+1$ in its two rows, with a single orientation $\sigma_j$ throughout
the run. Orientations are arbitrary subject to one rule: a district $i$
with $a_i = 0$ lies in a single row, so when both of its runs are
nonempty, $\sigma_{i-1}$ and $\sigma_i$ must place it in the same row.
Distinct admissible codings produce distinct partitions, with one
exception: when $a_i = a_{i+1} = 0$ and $\mu_{i-1} = \mu_{i+1} = 0$ (a
\emph{co-spanning pair}, horizontally splitting the same $s$ columns),
the two choices of $\sigma_i$ give the same partition. Consequently:
\begin{enumerate}[(a)]
  \item $|\Bcal_k^0(2 \times n)| = \Theta_k(n^{k-1})$;
  \item all interface parities are determined: $\mu_1 \equiv s$ and
    $\mu_{j+1} \equiv s - \mu_j \pmod 2$;
  \item $t(G[\xi_i]) = \hat t_{a_i}$ for every district, so
    \[
      \bar\pi(\xi) \;\propto\; \prod_{i=1}^{k} \hat t_{a_i}
      \;=\; \Theta_k\bigl(\lambda^{-\sum_j \mu_j}\bigr),
      \qquad \lambda = 2 + \sqrt3,
    \]
    and the projected measure concentrates exponentially on minimal
    interface widths.
\end{enumerate}
\end{proposition}

Informally, a balanced map of the strip is a sequence of solid blocks
joined by slanted interfaces. Its probability depends only on the
interface widths, and each additional cell of width divides the
probability by roughly $\lambda \approx 3.73$.

\begin{proof}
\emph{Supports are intervals.} If district $D$ meets columns $c < c'$
but not an intermediate column $d$, a path inside $D$ joining the two
sides would use a cell of column $d$, which $D$ does not have.

\emph{Overlap structure.} A column has two cells, so it meets at most
two districts, and two districts sharing a column each occupy one of its
rows there. Order the districts by the left endpoints of their supports.

\emph{The full columns of a district are consecutive.} Suppose district
$D$ has full columns $c < c'$ and some intermediate column is not full
for $D$ (a \emph{pocket}). Since supports are intervals, $D$ meets every
pocket column; the remaining pocket cells belong to districts whose
supports lie strictly inside the pocket, as they cannot contain the full
columns $c, c'$. Exactly as in Step 2 of
Proposition~\ref{prop:strip-classify}, connectivity of $D$ across the
pocket forces one entire row of the pocket into $D$, so the other pocket
row is a disjoint union of trapped districts, each of them a path of $s$
cells; the pocket width is then at least $s$ and
$|D| \ge 2 + 2 + s > s$, contradicting balance.

\emph{Runs and orientations.} Between two consecutive full blocks every
column is mixed, and by the support ordering the mixed pairs appear as
$(i, i+1), (i+1, i+2), \ldots$ in left-to-right order; the intermediate
districts have no full column, hence contain no vertical edge, hence lie
in a single row. Within one maximal run the orientation is constant: a
flip creates a maximal one-row segment of one district whose horizontal
neighbors on both sides and all vertical neighbors belong to the other
district of the run, stranding it from the rest of its district exactly
as in Step 4 of Proposition~\ref{prop:strip-classify}; for a
row-confined district a flip contradicts row-confinement outright.

\emph{Balance and the converse.} District $i$ has $2a_i$ cells in its
full block and one cell in each column of its adjacent runs, giving
$2a_i + \mu_{i-1} + \mu_i = s$, which determines the $\mu_i$
recursively. Conversely, every admissible coding produces a balanced
connected $k$-partition: each district is a full block with pendant
fringes, or a single-row path. The column word determines the partition,
and the word determines the coding except that flipping a $\sigma_i$ can
at most permute district names, which yields the same unordered
partition exactly in the co-spanning case. For the count,
$(a_1, \ldots, a_{k-1})$ ranges over the lattice points of a
$(k-1)$-dimensional polytope with side lengths $\Theta(s)$, and
orientation multiplicities lie in $[1, 2^{k-1}]$. The parity claims
follow from the balance equations read from the left end ($\mu_0 = 0$);
the tree counts from the pendant-path argument of
Proposition~\ref{prop:strip-law}; and the product form from
$\sum_i a_i = n - \sum_j \mu_j$ together with
$\hat t_a = \Theta(\lambda^{a})$.
\end{proof}

\done{Proposition~\ref{prop:strip-classify-k} is machine-verified as a
set equality (coded family versus brute-force enumeration) for
$k = 2, 3, 4$ at all admissible $n \le 12$, $k = 5$ at
$n \in \{5, 10\}$, and $k = 6$ at $n = 9$, together with
$t(G[\xi_i]) = \hat t_{a_i}$ for every district of every partition
(Section~\ref{sec:computation}, item 7). This completes
Theorem~\ref{conj:strip} step 1 in full, for exact balance.}

\begin{proposition}[Interior windows]
\label{prop:windows}
Let $\xi \in \Bcal_k^0(2 \times n)$ and let $\{D, D'\}$ be any edge of
$Q_\xi$. Then $G[D \cup D']$ is a $2 \times w$ ladder, $0 \le w \le s$,
with a pendant path attached at a prescribed corner on each side (the
cells the pair contributes to runs shared with other districts; for
non-consecutive $Q_\xi$-neighbors the union is a ladder with an extended
pendant, or a single path when $w = 0$), and
$t(G[D \cup D']) = \hat t_w$. Writing $p, q$ for the pendant lengths
($p + q + 2w = 2s$), the balanced connected $2$-partitions of the window
are exactly the staircase splits
\[
  A = (\text{left pendant}) \cup (\text{first } x \text{ ladder
  columns}) \cup (\text{a fringe of } f \text{ cells in one row}),
\]
up to complementation, where $p + 2x + f = s$, $x, f \ge 0$, and
$x + f \le w$. In particular pendants are never severed, there are at
most $2(w+1)$ balanced splits, and Lemma~\ref{lem:split}(a) applies
verbatim. A split with $f \in \{0, 1\}$ always exists, and for
$w \ge 1$
\[
\begin{gathered}
  \Pr\bigl[\text{a uniform spanning tree of } G[D \cup D'] \text{
  admits a balanced split}\bigr] \\
  \;\ge\; \frac{2 \hat t_x \hat t_y}{t_w} \;\ge\; \frac{1}{10},
  \qquad x + y \ge w - 1,
\end{gathered}
\]
while for $w = 0$ the window is a path and the probability is $1$.
\end{proposition}

\begin{proof}
The shape statement is read off the coding of
Proposition~\ref{prop:strip-classify-k}: the pair's full blocks and
shared run contribute full columns (the ladder), and the cells the pair
holds in runs shared with third districts are one-row segments attached
at a ladder corner by a single edge (pendant paths); pendant paths do
not change tree counts, so $t(G[D \cup D']) = \hat t_w$.

For the splits: a balanced part meeting a pendant contains its
attachment corner, since a connected set strictly inside a pendant is a
path of fewer than $s$ cells (the case $p = s$, where one part equals
the whole pendant, is the split with $x = f = 0$). The two parts have
interval column supports, and the analysis of
Proposition~\ref{prop:strip-classify} applied to the ladder portion
yields exactly the staircase family; counting $x \in \{0, \ldots, w\}$
and at most two fringe rows gives at most $2(w+1)$ splits.

Existence: set $f \equiv s - p \pmod 2$ with $f \in \{0, 1\}$, and
$x = (s - p - f)/2$. Then $x \ge 0$ because $p \le s$ (district $D$'s
balance equation), and $x + f \le w$ because $q \le s$. The parts have
tree counts $\hat t_x$ and $\hat t_y$ with $y = w - x - f$, so
$x + y \ge w - 1$, and Lemma~\ref{lem:split}(a) gives the bound
$2 \hat t_x \hat t_y / t_w$ (a vertical or near-vertical split has at
least two boundary edges). Finally, for $x, y \ge 1$ the closed form
$t_a = (\lambda^a - \lambda^{-a})/(2\sqrt3)$ gives
\[
  \frac{t_x t_y}{t_{x+y}}
  = \frac{\lambda^{x+y} - \lambda^{x-y} - \lambda^{y-x} +
      \lambda^{-x-y}}{2\sqrt3\,(\lambda^{x+y} - \lambda^{-x-y})}
  \;\ge\; \frac{1 - 2\lambda^{-2}}{2\sqrt3} \;>\; \frac15,
\]
and the ratio equals $1$ when $x = 0$ or $y = 0$; combined with
$t_{x+y} \ge t_{w-1} \ge t_w/4$ (from the recurrence), the probability
is at least $2 \cdot \frac15 \cdot \frac14 = \frac1{10}$.
\end{proof}

\done{Proposition~\ref{prop:windows} is machine-verified on $125$ window
configurations ($w \le 6$, $p, q \le 3$, all corner placements): the
staircase family equals the brute-force set of balanced splits, the
split identity is exact, and the minimum observed splittability
probability is $0.40$ against the proven bound $0.10$
(Section~\ref{sec:computation}, item 8). This completes
Theorem~\ref{conj:strip} step 2 in full.}

\begin{corollary}[Irreducibility on strips, all $k$]
\label{cor:strip-irred}
For every $k \ge 2$ with $k \mid 2n$, the projected
window-resampling chain on $\Bcal_k^0(2 \times n)$ is irreducible, with
diameter at most $2(k-1)$.
\end{corollary}

\begin{proof}
Sweep the interfaces left to right. Districts $1$ and $2$ are always
$Q_\xi$-adjacent: they share vertical edges when $\mu_1 \ge 1$ (which
covers every configuration with $a_1 = 0$, since then $\mu_1 = s$), and
abut through full columns when $\mu_1 = 0$. By
Proposition~\ref{prop:windows} the resample at interface $1$ realizes
every admissible $(a_1', \mu_1', \sigma_1')$; choose $\mu_1'$ minimal
for its parity ($0$ or $1$). Inductively, once interface $j-1$ has
minimal width $\mu_{j-1}' \le 1$, districts $j$ and $j+1$ are adjacent:
if $\mu_j \ge 1$ they share a run, and if $\mu_j = 0$ then
$2a_j = s - \mu_{j-1}' \ge s - 1 \ge 1$, so district $j$ has a full
block abutting district $j+1$'s columns. Setting $\mu_j'$ minimal is
admissible because $\mu_j' \le \mu_j$ (equal parities), whence
$2a_{j+1}' = s - \mu_j' - \mu_{j+1} \ge s - \mu_j - \mu_{j+1} \ge 0$.
After one sweep the interface widths equal the parity vector and the
$a_i$ are determined; the only remaining freedom is the orientation of
width-one runs, and one further resample per interface realizes any
orientation. Hence every state reaches a common canonical configuration
in at most $2(k-1)$ moves.
\end{proof}

\done{Corollary~\ref{cor:strip-irred} settles the strip case of the
irreducibility program for \emph{all} $k$ under \emph{exact} balance
(compare Conjecture~\ref{conj:irred}, which concerns rectangles). Both
heat-bath projected chain was also verified irreducible by machine at
every tested $(k, n)$.}

\begin{remark}[Tree-based kernels]
\label{rem:tree-kernels}
Practical redistricting chains propose window moves through spanning
trees, and two cautions emerged from our audit; both are reasons the
results above are stated for the window resampling kernel. First, the
idealized tree proposal (a uniform spanning tree of the window, then a
uniform balanced split edge) admits no closed-form Metropolis--Hastings
correction, because the reverse-proposal probability averages the
reciprocal number of split edges over all window trees. Second, the
implementable corrections condition the proposal tree to pass through
the current split, and any such kernel can move between two splits only
if some window tree contains exactly one boundary edge of each. Under
exact balance the splits of a window cross one another, and exhaustive
computation shows that no such tree exists at $s = 4, 6, 8$; the
projected tree-based chain is then frozen even though each of its
updates is exactly reversible (Section~\ref{sec:computation}, item 11).
Chains whose state carries the forest itself, such as the cycle walk and
the balanced up-down walk~\cite{deford2025cycle, akitaya2026bud}, do not
face this issue, because the tree is part of the state rather than an
auxiliary draw.
\end{remark}

\paragraph{Organization of the proof.} The proof of
Theorem~\ref{conj:strip} has four ingredients. (1) \emph{Structure}:
Proposition~\ref{prop:strip-classify-k}. (2) \emph{One interface}:
Proposition~\ref{prop:windows}, which also makes the heat-bath update
explicitly computable, with Corollary~\ref{cor:strip-irred} providing
irreducibility. (3) \emph{Many
interfaces}: Lemmas~\ref{lem:concentration} and~\ref{lem:coupling}. (4)
\emph{Assembly}: the proofs at the end of
Section~\ref{subsec:coupling}.

\paragraph{Exact spectral gaps.} The random-scan heat-bath projected
chain was diagonalized exactly on the coded state space:
\begin{center}
\begin{tabular}{cccc}
\hline
$k$ & $n$ & states & $\gamma(\bar P)$ \\
\hline
$3$ & $6$  & $9$   & $0.400$ \\
$3$ & $12$ & $33$  & $0.492$ \\
$3$ & $18$ & $73$  & $0.499$ \\
$3$ & $24$ & $129$ & $0.500$ \\
$4$ & $8$  & $25$  & $0.238$ \\
$4$ & $12$ & $77$  & $0.306$ \\
$4$ & $16$ & $177$ & $0.326$ \\
$5$ & $10$ & $64$  & $0.168$ \\
$5$ & $15$ & $288$ & $0.227$ \\
\hline
\end{tabular}
\end{center}
The gap \emph{increases} with $n$ toward $1/(k-1)$, the maximum
possible for a random-scan chain on $k-1$ coordinates, and the chain is
reversible to machine precision and irreducible at every tested
$(k, n)$.

\subsection{Proof of the strip gap}
\label{subsec:coupling}

In the coordinates of Proposition~\ref{prop:strip-classify-k} the
heat-bath chain admits a simple description. Pick one of the $k-1$
interfaces at random and redraw its width and orientation from the
correct conditional distribution, leaving everything else unchanged. The
conditional charges a factor of about $\lambda \approx 3.73$ per cell of
width, so it is heavily biased toward the narrowest admissible widths.
The next lemma makes this quantitative in two ways, and neither
statement cares what the rest of the partition looks like.

\begin{lemma}[Concentration and merging of the conditionals]
\label{lem:concentration}
Consider the heat-bath update at an interface whose two neighboring
interfaces currently have widths $u$ and $v$, and let
$\varepsilon \in \{0, 1\}$ be the parity of the updated width
(Proposition~\ref{prop:strip-classify-k}(b)).
\begin{enumerate}[(a)]
  \item \textbf{Concentration.} For every $u, v$ and every $r \ge 0$,
    the new width $m$ satisfies
    $\Pr[m \ge \varepsilon + 2r] \le 31\, \lambda^{-2r}$.
  \item \textbf{Merging.} If $s \ge 8$ and $u, v, u', v' \le s/4$
    (with the parities they must have), the conditionals given $(u, v)$
    and given $(u', v')$ differ in total variation by at most
    $40\, \lambda^{-s/4}$.
\end{enumerate}
\end{lemma}

\begin{proof}
Up to a factor not depending on $m$, the weight of width $m$ is
$\lambda^{-m} g(m)$ with $g(m) = h(a)\, h(a')\, \nu(m)$, where
$a = (s - u - m)/2$ and $a' = (s - m - v)/2$ are the two block widths,
$h(a) := \hat t_a / (A \lambda^a)$ with $A = 1/(2\sqrt3)$, and
$\nu(m) \in [1, 2]$ counts the admissible orientations. From the closed
form, $h(a) = 1 - \lambda^{-2a} \in [1 - \lambda^{-2}, 1)$ for
$a \ge 1$, and $h(0) = 2\sqrt3$. Hence $g(m) \in [c_-, c_+]$ with
$c_+ / c_- \le 2\,(2\sqrt3)^2 / (1 - \lambda^{-2})^2 \le 28$.

(a) The smallest admissible width is $\varepsilon$ whenever the support
is nonempty, so
\[
  \Pr[m \ge \varepsilon + 2r]
  \;\le\; \frac{c_+ \sum_{i \ge r} \lambda^{-(\varepsilon + 2i)}}
               {c_- \lambda^{-\varepsilon}}
  \;=\; \frac{c_+}{c_-} \cdot \frac{\lambda^{-2r}}{1 - \lambda^{-2}}
  \;\le\; 31\, \lambda^{-2r} .
\]
Truncating the support from above only helps, and the degenerate
point-mass cases are trivial.

(b) On widths $m \le s/2$ both conditionals are supported (the largest
admissible width is at least $3s/4$), both blocks satisfy
$a, a' \ge s/8 \ge 1$, so no orientation constraint is active and
$\nu$ is the same, while the factors $1 - \lambda^{-2a}$ of the two
conditionals differ by a multiplicative $1 \pm 3 \lambda^{-s/4}$. Widths
above $s/2$ carry mass at most $31 \lambda^{-s/2 + 2}$ under either
conditional, by (a). A ratio comparison on the common part plus the two
tails gives total variation at most $40\, \lambda^{-s/4}$.
\end{proof}

\begin{lemma}[Coupling of the interface chain]
\label{lem:coupling}
There are absolute constants $C_0, c > 0$ such that for every $k \ge 2$
and every $s \ge C_0 \log(k+1)$, the random-scan heat-bath interface
chain satisfies
\[
  t_{\mathrm{mix}}\bigl(P^{\mathrm{HB}}, 1/4\bigr)
  \;\le\; 2(k-1)\bigl(\lceil \log(k+1) \rceil + 4\bigr),
  \qquad
  \gamma\bigl(P^{\mathrm{HB}}\bigr) \;\ge\; \frac{c}{k \log(k+1)} .
\]
\end{lemma}

\begin{proof}
Run two copies $X$ and $Y$ from arbitrary states. At every step both
copies update the \emph{same} uniformly chosen interface, and the two
new (width, orientation) pairs are drawn from the maximal coupling of
the two conditionals. Let $M = \lfloor s/4 \rfloor$, let each of phase 1
and phase 2 last $T = (k-1)(\lceil \log(k+1) \rceil + 4)$ steps, and
consider three events.

First, in each phase every interface is selected at least once; missing
one is a coupon-collector event of probability at most
$(k-1) e^{-T/(k-1)} \le e^{-4}$ per phase.

Second, every update in either phase lands at width at most $M$ in both
copies. A single update fails with probability at most
$31 \lambda^{-M + 1}$ by Lemma~\ref{lem:concentration}(a), whatever the
neighbors are, so a union bound over at most $4T$ update-copies gives
$124\, T \lambda^{-M + 1}$. An interface at which the update holds (its
window is disconnected, which requires a neighboring interface of width
$s$) already has width $0 \le M$. Hence after phase 1 every interface
width in both copies is at most $M$, and this persists through phase 2.

Third, every phase-2 update draws the same pair in both copies. Given
all widths at most $M \le s/4$, no update holds (a hold needs a neighbor
of width $s$), and Lemma~\ref{lem:concentration}(b) bounds the failure
of each maximal coupling by $40 \lambda^{-s/4}$, so a union bound gives
$40\, T \lambda^{-s/4}$.

On the intersection of the three events, every phase-2 update leaves the
updated interface equal in the two copies; in particular, after the last
update of each interface the copies agree there, and that interface is
not modified again. Since every interface is updated during phase 2,
$X = Y$ after $2T$ steps. The total failure probability is at most
$2 e^{-4} + 124\, T \lambda^{-M+1} + 40\, T \lambda^{-s/4} \le 1/4$
once $s \ge C_0 \log(k+1)$ for an absolute $C_0$. The coupling
inequality gives the mixing bound, and
$t_{\mathrm{mix}}(1/4) \ge (t_{\mathrm{rel}} - 1) \log 2$ converts it
into the gap bound.
\end{proof}

\begin{proof}[Proof of Theorem~\ref{conj:strip}]
The chain $\bar P$ is exactly the random-scan heat-bath interface
chain, so Lemma~\ref{lem:coupling} bounds its mixing time, and
$t_{\mathrm{mix}}(1/4) \ge (t_{\mathrm{rel}} - 1) \log 2$ converts the
bound into the gap bound. For smaller $s$ the chain is a finite,
reversible, irreducible chain (Corollary~\ref{cor:strip-irred}), so its
gap is positive.
\end{proof}

\begin{proof}[Proof of Corollary~\ref{cor:strip-mixing}]
By Theorem~\ref{thm:conditional},
$t_{\mathrm{mix}}(P, \epsilon) = t_{\mathrm{mix}}(\bar P, \epsilon)$,
and $\bar P$ is exactly the random-scan heat-bath interface chain, so
Lemma~\ref{lem:coupling} gives the mixing bound directly; the coupling
bounds total variation itself, so no spectral argument and no laziness
are needed. The finitely many $n$ with $s < C_0 \log(k+1)$ contribute a
constant by Corollary~\ref{cor:strip-irred}, and $k = 2$ is
Remark~\ref{rem:k2-trivial}. For the cost per step, $P_F$ runs Wilson's
algorithm in every district and $P_M$ resamples one window, with the
split distribution computed from Proposition~\ref{prop:windows} and the
ladder recurrence.
\end{proof}

\done{Lemmas~\ref{lem:concentration} and~\ref{lem:coupling},
Theorem~\ref{conj:strip}, and Corollary~\ref{cor:strip-mixing} are
proved in full. Both estimates of Lemma~\ref{lem:concentration} are
machine-verified with wide margins: the tail constant has a $33\%$
cushion over all neighbor widths at $s \le 40$, and the measured merging
total variation is three or more orders of magnitude below the stated
bound (Section~\ref{sec:computation}, item 10).}

\statusnote{Scope: the strip program is complete for exact balance; the
$\varepsilon$-balanced variant of the classification is routine but not
written. Torpidity examples of~\cite{charikar2023complexity} live on
grid-\emph{like} planar families and do not touch strips; any claim
beyond constant-height strips must engage the geometry
of~\cite{cannon2024balanced} seriously. Square grids are stated as the
open problem, not promised.}

%═══════════════════════════════════════════════════════════════════════════════
\section{Machine verification}
\label{sec:computation}
%═══════════════════════════════════════════════════════════════════════════════

All computations are exhaustive (no sampling) and were performed in exact
integer or rational arithmetic. Scripts are available from the author.

\begin{enumerate}[1.]
  \item \textbf{Torsor axioms (Theorem~\ref{thm:torsor}).} For the
    $2 \times 3$ grid with $S$ a pair of opposite corners ($|\Fcal_S| =
    21$) and with a root triple ($13$), the $3 \times 3$ grid with opposite
    corners ($288$), and $K_4$ with an adjacent pair ($8$; here $G_S$ has
    parallel edges): every generator acts as a permutation of
    $\Tcal(G_S)$; the chip-firing relations hold, so the action factors
    through $\Pick{G}$; generators commute; and the orbit of a single tree
    has full size. Freeness follows from transitivity and
    $|\Pick{G}| = |\Fcal_S(G)|$.
  \item \textbf{Matrix-forest consistency.} For \emph{every} root set of
    every test graph (triangle, $C_4$, $K_4$, $2\times 3$, $3 \times 3$;
    $162$ root sets in total): $\det \Delta^{(S)} = |\Fcal_S(G)| = t(G_S)$,
    and the Smith normal forms of $\Delta^{(S)}$ and of the reduced
    Laplacian of $G_S$ agree.
  \item \textbf{Fiber decomposition (Theorem~\ref{thm:fiber-decomp}).}
    Verified by enumeration on five partitions of the $2 \times 3$ and
    $3 \times 3$ grids ($\tau = 3, 4, 8, 45, 32$) and on a $6 \times 6$
    quadrant instance ($\tau = 192^4 \cdot 108$).
  \item \textbf{No algebraic shortcut for balance.} On the $2 \times 3$
    grid with $S = $ opposite corners: of the $9$ partition fibers in
    $\Pick{G} \cong$ ($21$ elements), none of the fibers of size $> 1$ is a
    coset of a subgroup, and the union of the three balanced fibers is not
    a coset. Moreover no group homomorphism $\PicO{G} \to \Pick{G}$ can
    exist in general ($|{\Pick{}}| \nmid t(G)$ for $162$ of the tested root
    sets, including the triangle with $k = 2$: $2 \nmid 3$).
  \item \textbf{Irreducibility evidence.} On the $3 \times 3$ grid,
    $k = 2$, no balance constraint: the boundary-restricted single-swap
    chain on (tree, marked-edge) states is irreducible: all $1{,}536$
    states form one class covering all $53$ connected $2$-partitions. With
    exact balance ($k = 3$, districts of size $3$): the state graph has two
    isolated states, but the main class covers all $10$ balanced
    partitions, consistent with Conjecture~\ref{conj:irred} at the
    partition level and a warning at the lifted level.
  \item \textbf{Balanced partition counts on strips.} Exhaustive counts
    of $|\Bcal_k^0(2 \times n)|$: for $k = 2$, exactly $n$ for all
    $n \le 12$, matching Proposition~\ref{prop:strip-classify}; for
    $k = 3$: $3, 9, 19, 33$ at $n = 3, 6, 9, 12$; for $k = 4$:
    $5, 11, 25, 45, 77$ at $n = 4, 6, 8, 10, 12$, consistent with
    $\Theta(n^{k-1})$ growth (constant second differences at $k = 3$,
    cubic growth at $k = 4$).
  \item \textbf{General classification
    (Proposition~\ref{prop:strip-classify-k}).} The coded family equals
    the brute-force set of balanced connected $k$-partitions, as sets of
    partitions, for $k = 2, 3, 4$ at all admissible $n \le 12$, $k = 5$
    at $n \in \{5, 10\}$ ($8$ and $64$ partitions), and $k = 6$ at
    $n = 9$ ($41$ partitions); and $t(G[\xi_i]) = \hat t_{a_i}$ holds for
    every district of every coded partition (matrix-tree check). The
    balanced-forest multiplicities on $(k, n) = (3, 6)$ equal
    $\prod_i \hat t_{a_i}$ for all $9$ partitions, confirming the
    projected law.
  \item \textbf{Interior windows (Proposition~\ref{prop:windows}).} On
    $125$ pendant-ladder configurations ($w \le 6$, $p, q \le 3$, all
    corner placements): the staircase family equals the brute-force set
    of balanced splits, the split-count identity is exact for every
    split, $t(H) = \hat t_w$, and the minimum splittability probability
    is $0.400$ (proven bound $1/10$). The pendant-free rectangular cases
    ($w = 4, 6, 8$; $56$, $780$, $10{,}864$ trees) give splittability
    $0.786$, $0.764$, $0.763$ against the limit $1/\sqrt3 \approx
    0.577$.
  \item \textbf{Projected spectral gaps.} Exact diagonalization of
    the heat-bath chain at the $(k, n)$ pairs of the table in
    Section~\ref{subsec:strips}: reversible to machine precision,
    irreducible everywhere, and $\gamma(\bar P) \to 1/(k-1)$ as $n$
    grows.
  \item \textbf{Coupling estimates
    (Lemma~\ref{lem:concentration}).} Exact conditionals at all
    admissible neighbor widths for $s \le 40$: the tail bound holds with
    maximum ratio $0.67$ of the stated constant, and the merging total
    variation at $s = 16, 24, 32, 40$ is smaller than $\lambda^{-s/4}$
    by three or more orders of magnitude.
  \item \textbf{Tree-based kernel freeze
    (Remark~\ref{rem:tree-kernels}).} The through-the-current-split
    merge--split kernel satisfies detailed balance to machine precision,
    yet its off-diagonal support is empty at every tested size: at
    $(k, n)$ with $s = 4, 6, 8$ the projected chain has spectral gap
    exactly $0$ and is reducible.
\end{enumerate}

\done{All five computations were run to completion in July 2026.}
\statusnote{Before submission: re-run the suite from a clean repository,
commit the scripts (currently in a scratch directory), and add a $4 \times
4$, $k = 2$ instance of item 5, which is within reach ($t(4\times 4)
= 100{,}352$ trees; restrict to balanced marked states).}

%═══════════════════════════════════════════════════════════════════════════════
\section{Discussion}
%═══════════════════════════════════════════════════════════════════════════════

The division of labor proposed here is strict: algebra handles everything
that is exactly solvable (uniformity on fibers, torsor navigation,
within-fiber waste, which Lemma~\ref{lem:collapse} eliminates rather than
mitigates), and the remaining difficulty is a spectral gap estimate for an
explicit chain on an explicit state space, which
Section~\ref{sec:targets} supplies on strips
(Theorem~\ref{conj:strip}). We emphasize what the collapse lemma does
\emph{not} do. It does not make the projected chain fast; torpidity
examples~\cite{charikar2023complexity} project too. It makes the projected
chain the \emph{whole} problem, which no existing analysis of a balanced
chain achieves. With the strip program complete, genuine two-dimensional
grids are the central open problem.

Two further directions are natural. First, the multi-sink torsor suggests
global $P_M$ moves: resample the entire forest conditioned on a root set
via Wilson on $G_S$, with roots re-drawn from a suitable distribution; the
acceptance ratio is computable by Theorem~\ref{thm:fiber-decomp}, and such
moves may shortcut the interface-by-interface diffusion on strips. Second,
the merge--split ledger of Xu's partition divisors (the minimal effective
divisor $D^\xi$ with $\deg D^\xi = |\partial \xi|$, and the merge formula
of \cite[Cor.~3]{xu2011partitions}; see also the resolution theorems
of~\cite{manjunath2015minimal, mohammadi2014divisors}) provides an exact
algebraic bookkeeping of merge--split paths that may serve canonical-path
constructions on general planar graphs.

\missing{Decide, after Theorem~\ref{conj:strip} is resolved, whether the
global-move variant belongs in this paper or the next one.}

%═══════════════════════════════════════════════════════════════════════════════
\appendix
\section{Status}
\label{app:status}
%═══════════════════════════════════════════════════════════════════════════════

\begin{center}
\begin{tabular}{p{7.6cm} p{6.8cm}}
\hline
\textbf{Item} & \textbf{Status} \\
\hline
Thm~\ref{thm:torsor} (multi-sink torsor, (i)--(iii)) &
  \textcolor{donecol}{Proved in full; machine-verified.} \\
Rem.~\ref{rem:dhar} (multi-sink Dhar) &
  \textcolor{donecol}{Immediate from Thm~\ref{thm:torsor}; verified.} \\
Cor.~\ref{cor:sampling}(a) (Wilson exact sampling) &
  \textcolor{donecol}{Proved (citation to Wilson '96).} \\
Cor.~\ref{cor:sampling}(b) (navigation) &
  \textcolor{donecol}{Proved as stated;}
  \textcolor{notecol}{torsor discrete-log deliberately avoided.} \\
Thm~\ref{thm:fiber-decomp} (fiber decomposition) &
  \textcolor{donecol}{Proved in full; verified numerically.} \\
Prop.~\ref{prop:resampler} (exact fiber resampler) &
  \textcolor{donecol}{Proved in full.} \\
Lem.~\ref{lem:collapse} (fiber collapse, (a)--(d)) &
  \textcolor{donecol}{Proved in full.} \\
Thm~\ref{thm:conditional} (conditional mixing theorem) &
  \textcolor{donecol}{Proved in full (conditional on projected gap).} \\
$P_M$ specification, MH correction details (\S\ref{subsec:pm}) &
  \textcolor{misscol}{Write out the exact proposal/acceptance formulas;
  routine but must be in the paper.} \\
Conj.~\ref{conj:irred} step 1 ($k=3$ rectangles) &
  \textcolor{misscol}{Open; adapt \cite[App.~B]{akitaya2026bud};
  first priority.} \\
Conj.~\ref{conj:irred} steps 2--3 (general $k$) &
  \textcolor{misscol}{Open; competition risk (BUD authors).} \\
BUD $\to$ merge--split transfer ($k=2$ base case) &
  \textcolor{misscol}{Believed routine via \cite[Lem.~A.3--A.5]{akitaya2026bud};
  write it.} \\
Thm~\ref{conj:strip} step 1, $k = 2$: classification
(Prop.~\ref{prop:strip-classify}, exactly $n$ staircases) and stationary
law (Prop.~\ref{prop:strip-law}) &
  \textcolor{donecol}{Proved in full; counts machine-verified to
  $n = 12$.} \\
$k = 2$ triviality of merge--split (Rem.~\ref{rem:k2-trivial});
Theorem~\ref{conj:strip} restated for $k \ge 3$ &
  \textcolor{donecol}{Done.} \\
Thm~\ref{conj:strip} step 1, $k \ge 3$: general classification,
parities, product-form stationary law
(Prop.~\ref{prop:strip-classify-k}) &
  \textcolor{donecol}{Proved in full; set equality machine-verified for
  $k \le 6$.} \\
Strip irreducibility for all $k$, exact balance, diameter $\le 2(k-1)$
(Cor.~\ref{cor:strip-irred}) &
  \textcolor{donecol}{Proved in full.} \\
$P_M$ implementability audit: tree-based MH kernels frozen under
exact balance; theorem chain switched to window heat-bath
(Rem.~\ref{rem:tree-kernels}) &
  \textcolor{donecol}{Done; freeze machine-verified at $s = 4, 6, 8$.} \\
Thm~\ref{conj:strip} step 2 on rectangular windows: split identity,
constant split probability, $\mathrm{poly}$-accurate proposal
(Lem.~\ref{lem:split}) &
  \textcolor{donecol}{Proved in full; machine-verified $w = 4, 6, 8$.} \\
Thm~\ref{conj:strip} step 2 on interior windows: shapes, split
classification, splittability $\ge 1/10$
(Prop.~\ref{prop:windows}) &
  \textcolor{donecol}{Proved in full; verified on $125$
  configurations.} \\
Thm~\ref{conj:strip} step 3: concentration + coupling
(Lems~\ref{lem:concentration}, \ref{lem:coupling}) &
  \textcolor{donecol}{Proved in full; estimates machine-verified with
  wide margins.} \\
Thm~\ref{conj:strip} and Cor.~\ref{cor:strip-mixing}: unconditional
strip mixing, $O(k \log k)$ steps, near-linear total time &
  \textcolor{donecol}{Proved in full: the headline theorem.} \\
Exact gap diagonalization ($k = 3, 4, 5$; $n \le 24$) &
  \textcolor{donecol}{Done; table in \S\ref{subsec:strips}.} \\
$4 \times 4$ balanced irreducibility computation &
  \textcolor{misscol}{Not yet run; feasible.} \\
Scripts committed to repository &
  \textcolor{misscol}{Currently in session scratch directory.} \\
Related-work completeness pass (incl.\ any 2026 preprints) &
  \textcolor{misscol}{Re-sweep arXiv before submission.} \\
Remove inline status boxes and the WORKING DRAFT banner &
  \textcolor{misscol}{Before submission.} \\
\hline
\end{tabular}
\end{center}

\paragraph{Reading order for finishing the paper.}
(1) Theorem~\ref{conj:strip} step 2 is the pivot: if the one-interface
split-probability bound works out on strips, everything else is assembly
and the paper's headline becomes Corollary~\ref{cor:strip-mixing}
unconditionally for $k = 2$; submit to SODA/STOC. (2) If step 2 stalls,
the fallback paper is Sections~\ref{sec:torsor}--\ref{sec:collapse} plus
Conjecture~\ref{conj:irred} step 1 proved (irreducibility, $k = 3$
rectangles, exact balance): a coherent algorithms paper for ESA/SODA.
(3) The infrastructure alone (Sections~\ref{sec:torsor}--\ref{sec:fiber})
is the minimal publishable unit, as a combinatorics note.

\begin{thebibliography}{99}

\bibitem{akitaya2026bud}
H.~A. Akitaya, S.~Cannon, G.~Herschlag, G.~Schoenbach, K.~Tapp, and
  J.~Tucker-Foltz.
\newblock The balanced up-down walk.
\newblock \emph{arXiv preprint arXiv:2602.11993}, 2026.

\bibitem{anari2021logconcave}
N.~Anari, K.~Liu, S.~{Oveis Gharan}, C.~Vinzant, and T.-D. Vuong.
\newblock Log-concave polynomials {IV}: approximate exchange, tight mixing
  times, and near-optimal sampling of forests.
\newblock In \emph{Proceedings of the 53rd Annual ACM SIGACT Symposium on
  Theory of Computing (STOC '21)}, pages 408--420, 2021.

\bibitem{autry2020multiscale}
E.~A. Autry, D.~Carter, G.~Herschlag, Z.~Hunter, and J.~C. Mattingly.
\newblock Multi-scale merge-split {M}arkov chain {M}onte {C}arlo for
  redistricting.
\newblock \emph{arXiv preprint arXiv:2008.08054}, 2020.

\bibitem{baker2018bernardi}
M.~Baker and Y.~Wang.
\newblock The {B}ernardi process and torsor structures on spanning trees.
\newblock \emph{International Mathematics Research Notices},
  2018(16):5120--5147, 2018.

\bibitem{cannon2024balanced}
S.~Cannon, W.~Pegden, and J.~Tucker-Foltz.
\newblock Sampling balanced forests of grids in polynomial time.
\newblock In \emph{Proceedings of the 56th Annual ACM Symposium on Theory of
  Computing (STOC '24)}, pages 1676--1687, 2024.

\bibitem{chan2015rotor}
M.~Chan, T.~Church, and J.~A. Grochow.
\newblock Rotor-routing and spanning trees on planar graphs.
\newblock \emph{International Mathematics Research Notices},
  2015(11):3225--3244, 2015.

\bibitem{charikar2023complexity}
M.~Charikar, P.~Liu, T.~Liu, and T.-D. Vuong.
\newblock On the complexity of sampling redistricting plans.
\newblock \emph{arXiv preprint arXiv:2206.04883}, 2023.

\bibitem{chebotarev2006forest}
P.~Chebotarev and E.~Shamis.
\newblock Matrix-forest theorems.
\newblock \emph{arXiv preprint arXiv:math/0602575}, 2006.

\bibitem{deford2021recom}
D.~DeFord, M.~Duchin, and J.~Solomon.
\newblock Recombination: a family of {M}arkov chains for redistricting.
\newblock \emph{Harvard Data Science Review}, 3(1), 2021.

\bibitem{deford2025cycle}
D.~R. DeFord, G.~Herschlag, and J.~C. Mattingly.
\newblock A cycle walk for sampling measures on spanning forests for
  redistricting.
\newblock \emph{arXiv preprint arXiv:2509.08629}, 2025.

\bibitem{gallagher2025balanced}
M.~Gallagher and K.~Tapp.
\newblock Balanced spanning trees of the 2-by-$n$ grid.
\newblock \emph{arXiv preprint arXiv:2508.10614}, 2025.

\bibitem{ganguly2023rotor}
A.~Ganguly and A.~McDonough.
\newblock Rotor-routing induces the only consistent sandpile torsor structure
  on plane graphs.
\newblock \emph{Forum of Mathematics, Sigma}, 11:e78, 2023.

\bibitem{holroyd2008chip}
A.~E. Holroyd, L.~Levine, K.~M\'esz\'aros, Y.~Peres, J.~Propp, and D.~B.
  Wilson.
\newblock Chip-firing and rotor-routing on directed graphs.
\newblock In \emph{In and Out of Equilibrium 2}, volume~60 of \emph{Progress
  in Probability}, pages 331--364. Birkh\"auser, Basel, 2008.

\bibitem{madras2002decomposition}
N.~Madras and D.~Randall.
\newblock {M}arkov chain decomposition for convergence rate analysis.
\newblock \emph{Annals of Applied Probability}, 12(2):581--606, 2002.

\bibitem{manjunath2015minimal}
M.~Manjunath, F.-O. Schreyer, and J.~Wilmes.
\newblock Minimal free resolutions of the {$G$}-parking function ideal and
  the toppling ideal.
\newblock \emph{Transactions of the American Mathematical Society},
  367(4):2853--2874, 2015.

\bibitem{mew}
A.~McWhorter and D.~DeFord.
\newblock The {M}arked {E}dge {W}alk: a novel {MCMC} algorithm for sampling
  of graph partitions.
\newblock \emph{arXiv preprint arXiv:2510.17714}, 2025.

\bibitem{mohammadi2014divisors}
F.~Mohammadi and F.~Shokrieh.
\newblock Divisors on graphs, connected flags, and syzygies.
\newblock \emph{International Mathematics Research Notices},
  2014(24):6839--6905, 2014.

\bibitem{shokrieh2021torsor}
F.~Shokrieh and C.~Wright.
\newblock Torsor structures on spanning trees.
\newblock \emph{arXiv preprint arXiv:2103.10370}, 2021.

\bibitem{wilson1996generating}
D.~B. Wilson.
\newblock Generating random spanning trees more quickly than the cover time.
\newblock In \emph{Proceedings of the 28th Annual ACM Symposium on Theory of
  Computing (STOC '96)}, pages 296--303, 1996.

\bibitem{xu2011partitions}
T.~Xu.
\newblock \emph{Graph Partitions and Free Resolutions of Toppling Ideals}.
\newblock Bachelor's thesis, Reed College, 2011.

\end{thebibliography}

\end{document}
